% ==============================================================================
% METADATOS DE AUDITORÍA EDITORIAL Y RIGOR CIENTÍFICO (PROTOCOLO AGENTS.md)
% ==============================================================================
% Capítulo: 14
% Título: Ética, Plagio y Propiedad Intelectual en la Universidad del Siglo XXI: Integridad Académica, Derechos de Autor y el Desafío de la Inteligencia Artificial Generativa
% Grupo de Origen: Grupo 9 (Asignatura de Inteligencia Artificial Aplicada, UNMSM)
% Archivo Fuente Analizado: Capitulo_14_Etica_y_Plagio_Mejorado.pdf (27 páginas, 11,151 palabras)
% Fuentes Primarias Cotejadas y Verificadas (referencias/):
%   1. International Center for Academic Integrity [ICAI]. (2021). The fundamental values of academic integrity (3.ª ed.).
%   2. Perú. (1996, 24 de abril). Decreto Legislativo N.º 822 Ley sobre el Derecho de Autor. WIPO Lex.
%   3. American Psychological Association [APA]. (2022). Plagiarism and academic integrity guidelines.
%   4. De Jesús G., M. I. (2016). Entre la ética en la investigación y la propiedad intelectual. Actualidad Contable FACES, 19(32), 40–67.
%   5. Porto Castro, A. M. (2024). El plagio en la universidad: Una mirada pedagógica desde el rigor de los datos. Editorial Redipe.
%   6. Corona Domínguez, B., & Alonso Nuñez, M. C. (2026). Plagio académico y herramientas de inteligencia artificial generativa. Ciencia y Reflexión, 5(1), 797–825.
%   7. Puche Villalobos, D. J. (2025). La inteligencia artificial y el fraude académico en el contexto universitario. Revista Digital de Investigación y Postgrado, 6(11), 73–93.
%   8. Miao, F., & Holmes, W. (2023). Guidance for generative AI in education and research. UNESCO Publishing.
% Cláusula NotebookLM: Grounding estricto al 100%. Cero datos o métricas inventadas.
% Parámetro de Extensión: >= 25 páginas compiladas en formato oficial UNMSM (A4, 12pt).
% Decisiones del Usuario: Volcado integral de la investigación monográfica del Grupo 9 preservando todas sus tablas, casos, taxonomías y aparato empírico.
% Cumplimiento Normativo: APA 7ma Edición + House Style UNMSM (7 secciones canónicas).
% Estado de Auditoría: CONFORME / REVISADO POR CAPITULO-AUDITOR
% Fecha de Auditoría: 2026-11-20
% ==============================================================================

\section[Ética, Plagio y Propiedad Intelectual]{Ética, Plagio y Propiedad Intelectual en la Universidad del Siglo XXI: Integridad Académica, Derechos de Autor y el Desafío de la Inteligencia Artificial Generativa}
\markboth{CAPÍTULO 14\quad ÉTICA Y PROPIEDAD INTELECTUAL}{}

\begin{cajabeige}
\textbf{Pregunta Central del Capítulo:} \textquestiondown Cómo puede la universidad del siglo XXI aprovechar las herramientas digitales y la inteligencia artificial generativa, asegurando al mismo tiempo que los trabajos académicos expresen un aprendizaje genuino, una atribución rigurosa de las fuentes y una responsabilidad ética inquebrantable sobre su contenido?\\[0.3cm]
\textbf{Objetivo de Aprendizaje:} El lector comprenderá de manera exhaustiva los fundamentos axiológicos de la integridad académica y el derecho de autor según la legislación nacional e internacional; analizará críticamente las morfologías del plagio, las autorías ficticias y los sesgos metodológicos del software de detección; evaluará el impacto epistemológico de la inteligencia artificial generativa distinguiendo el andamiaje formativo de la sustitución fraudulenta; y diseñará instrumentos institucionales de gobernanza, rúbricas analíticas de proceso, cronogramas de trazabilidad y protocolos de debido proceso aplicables en la docencia e investigación universitaria.
\end{cajabeige}

\textoside{
La integridad académica en la era digital no constituye una mera postura reactiva de vigilancia o sanción disciplinaria, sino un compromiso axiológico activo fundamentado en seis valores universales: honestidad, confianza, justicia, respeto, responsabilidad y valentía. En las ciencias exactas y de la computación, cada investigación debe erigirse como un mapa epistemológico transparente donde se delimite con precisión milimétrica qué conceptos, datos, algoritmos o códigos proceden de fuentes ajenas y qué síntesis, demostración o innovación analítica constituye el aporte genuino de quien suscribe la obra.
}{CONCEPTO CLAVE\\[0.15cm]Valores de Integridad (ICAI)\\[0.1cm]Honestidad $\cdot$ Confianza $\cdot$ Justicia $\cdot$ Respeto $\cdot$ Responsabilidad $\cdot$ Valentía}

\vspace{0.3cm}

El acceso ubicuo a redes informáticas de alta velocidad, repositorios globales de código abierto y plataformas indexadas de literatura científica ha democratizado, a una escala sin precedentes históricos, la circulación y el consumo del conocimiento humano. Sin embargo, esta inmediatez técnica conlleva una vulnerabilidad epistemológica inherente: disponer de un archivo en formato PDF, consultar un conjunto masivo de datos (\textit{dataset}) o descargar un bloque de código ejecutable desde un repositorio digital no equivale en absoluto a comprender los principios subyacentes, ni autoriza en ninguna circunstancia a presentarlos ante la comunidad académica como un fruto del propio intelecto.

En disciplinas formales y aplicadas como la Computación Científica, las Matemáticas Puras y la Estadística, la autenticidad del trabajo intelectual reviste un valor fundacional. Cada manuscrito, algoritmo o demostración formal debe operar con absoluta transparencia técnica: quien lee el trabajo debe ser capaz de reconstruir el itinerario cognitivo completo, verificando de manera inequívoca qué axiomas, librerías, arquitecturas o conclusiones pertenecen al acervo colectivo de la ciencia y qué deducción original, refinamiento algorítmico o modelado predictivo aporta quien firma la investigación.

Esta demarcación, constitutiva del método científico moderno, enfrenta hoy un punto de inflexión sin precedentes. La consolidación masiva de herramientas fundamentadas en modelos generativos de lenguaje (\textit{Large Language Models}, LLMs) y sistemas de programación asistida permite sintetizar prosa coherente, traducir textos especializados, inferir funciones algorítmicas complejas y estructurar revisiones bibliográficas en cuestión de segundos. En este nuevo ecosistema cognitivo, la interrogante deontológica de la comunidad científica ya no se circunscribe al dilema elemental de \textquestiondown hubo copia textual o no?, sino que se traslada hacia una demanda de rendición de cuentas epistemológica: \textquestiondown posee el autor la capacidad crítica de explicar, sustentar, depurar y responder éticamente por cada ecuación, línea de código y afirmación teórica que suscribe bajo su firma?

El presente capítulo monográfico, formulado a partir de la exhaustiva investigación desarrollada por el equipo de la Escuela Profesional de Computación Científica de la Facultad de Ciencias Matemáticas de la Universidad Nacional Mayor de San Marcos (UNMSM) ---integrado por Palacios Carrasco Roberto (líder), Avalos Palomino Ismael, Noel Vargas Alexander Henry, Cajas Aquino Luis Alvaro y Lavado Choque Andrew Xavier---, ofrece una sistematización rigurosa de este panorama multidimensional. A través de un enfoque articulado que integra doctrina jurídica comparada, taxonomías de fraude, análisis empírico y gobernanza institucional, se dota a la comunidad universitaria de un marco teórico y operativo robusto para liderar la investigación científica con excelencia, integridad y responsabilidad ética.

\subsection{Integridad Académica, Propiedad Intelectual y Regímenes Jurídicos}

\subsubsection{Fundamentos Epistemológicos y Axiológicos de la Integridad Científica}

La integridad académica no puede concebirse de forma simplista como la ausencia pasiva de sanciones disciplinarias o la mera evasión de mecanismos automáticos de detección de similitud. Muy por el contrario, constituye un principio afirmativo y una cultura institucional proactiva orientada a consolidar la veracidad y reproducibilidad del conocimiento científico. Como postula el \textit{International Center for Academic Integrity} (ICAI, 2021), la integridad se edifica sobre la articulación armónica de seis valores fundamentales: honestidad, confianza, justicia, respeto, responsabilidad y valentía.

La significación práctica de estos valores no reside en enunciados abstractos o declamaciones retóricas, sino en su capacidad para normar decisiones metodológicas cotidianas en el laboratorio computacional, la biblioteca especializada y la redacción de informes técnicos. La \textbf{honestidad} impone el deber imperativo de no atribuirse méritos ajenos ni alterar observaciones experimentales; la \textbf{confianza} se sustenta en la transparencia de datos y métodos que permiten a la comunidad replicar los hallazgos; la \textbf{justicia} exige un reconocimiento escrupuloso y equitativo de las contribuciones intelectuales de todos los colaboradores; el \textbf{respeto} honra el esfuerzo de los autores precedentes mediante la citación exhaustiva; la \textbf{responsabilidad} asume plenamente las implicancias teóricas y sociales de lo publicado; y la \textbf{valentía} demanda defender la verdad científica, corregir públicamente los errores propios y negarse a prácticas fraudulentas, aun ante presiones jerárquicas o exigencias de productividad académica apresurada.

\begin{table}[htbp]
\centering
\small
\renewcommand{\arraystretch}{1.25}
\textbf{Tabla 14.1}\\[0.15cm]
\textit{Operacionalización de los Seis Valores Fundamentales del ICAI en la Investigación Universitaria}\\[0.25cm]
\begin{tabularx}{\textwidth}{p{2.6cm}X p{5.5cm}}
\toprule
\textbf{Valor Fundamental} & \textbf{Definición Operativa en el Ámbito Académico} & \textbf{Decisión Concreta y Demostrable en el Ciclo de Trabajo} \\
\midrule
\textbf{Honestidad} & Presentar lo propio como genuinamente propio y lo ajeno con atribución rigurosa e inequívoca. & Declarar en notas al pie o en el repositorio el uso de fragmentos de código externo y librerías auxiliares. \\
\textbf{Confianza} & Proveer evidencia empírica suficiente y transparente que permita la auditoría y replicación pública. & Publicar el cuaderno de código fuente (\textit{Jupyter notebook}) con las semillas aleatorias e instrucciones de entorno. \\
\textbf{Justicia} & Reconocer las autorías de acuerdo con la magnitud del aporte intelectual real, evitando privilegios espurios. & Aplicar la Taxonomía CRediT para asignar formalmente los roles de cada integrante del equipo investigador. \\
\textbf{Respeto} & Valorar el tiempo y la producción intelectual de los pares mediante la lectura profunda y la cita debida. & Leer integralmente las fuentes primarias antes de citarlas, evitando atribuirles tesis no formuladas en el original. \\
\textbf{Responsabilidad} & Asumir la paternidad intelectual y las consecuencias epistemológicas y éticas de todo lo firmado. & Verificar personalmente cada referencia bibliográfica y dato factual sugerido por asistentes de IA generativa. \\
\textbf{Valentía} & Sostener la rectitud metodológica y ética incluso cuando implique costos académicos o incomodidad social. & Reportar de inmediato al docente o comité un sesgo muestral inadvertido o retractar un resultado no replicable. \\
\bottomrule
\end{tabularx}
\vspace{0.15cm}
\raggedright \scriptsize \textbf{Nota.} Adaptado de \textit{The Fundamental Values of Academic Integrity} (3.ª ed.), por International Center for Academic Integrity [ICAI], 2021.
\end{table}

La adhesión a estos principios adquiere una trascendencia singular en la \textbf{firma colectiva} de proyectos de investigación. En disciplinas de modelado matemático e informática aplicada, los proyectos suelen ejecutarse mediante equipos multidisciplinarios. Es una práctica deformada y éticamente inaceptable fragmentar la responsabilidad aduciendo desconocimiento de una sección redactada por un compañero. La firma solidaria de un manuscrito o reporte técnico confiere autoría moral y legal a todos los firmantes sobre la totalidad del documento; por consiguiente, cada integrante del equipo asume el deber ineludible de comprender, auditar y estar en condiciones de defender oralmente cada afirmación, demostración y algoritmo integrado en el entregable final.

\subsubsection{Dualidad Jurídica: Normativa Ética Académica frente al Derecho de Autor}

Uno de los equívocos más persistentes en las comunidades universitarias radica en homologar mecánicamente las normas deontológicas de la academia con las disposiciones punitivas del derecho de autor. Si bien ambas esferas convergen en la protección del trabajo intelectual, operan bajo naturalezas jurídicas, alcances epistémicos y mecanismos sancionatorios marcadamente disímiles.

El \textbf{derecho de autor} constituye una rama especializada del ordenamiento jurídico positivo cuya finalidad primaria es tutelar la forma original mediante la cual las ideas se expresan exteriormente (creaciones literarias, artísticas, científicas, programas computacionales), así como regular los derechos económicos derivados de su explotación. Por expresa disposición de los tratados internacionales y las leyes nacionales, el derecho de autor \textit{no protege las ideas, los métodos matemáticos, los conceptos científicos abstractos ni los algoritmos en sí mismos}, sino exclusivamente la formulación discursiva o el código concreto en el que quedan plasmados.

En contraste, la \textbf{ética académica} constituye un régimen normativo interno fundado en la confianza comunitaria y el avance acumulativo del saber. La ética universitaria sí tutela rigurosamente las ideas, las hipótesis metodológicas, los descubrimientos teóricos y los razonamientos algorítmicos. En consecuencia, una conducta puede ser perfectamente lícita ante los tribunales de justicia y, de forma simultánea, configurar una infracción ética gravísima merecedora de la expulsión académica o la anulación de un título universitario.

Desde la perspectiva de la filosofía del derecho y la epistemología jurídica, la justificación de la propiedad intelectual se remonta a la teoría del trabajo de John Locke (1690), según la cual el creador adquiere un derecho sobre aquello que produce al imprimir su esfuerzo intelectivo. No obstante, teóricos contemporáneos como Hettinger (1989) y Drahos (2016) advierten que las ideas no son bienes rivales ni excluyentes, por lo que el monopolio legal del autor debe balancearse cuidadosamente con el bien público del avance del conocimiento. Como demuestran Landes y Posner (2003) y Netanel (2008), la hipertrofia de los derechos de autor puede encarecer artificialmente la investigación y asfixiar la producción cultural colectiva, razón por la cual los tratados internacionales (Goldstein \& Hugenholtz, 2019; Ricketson \& Ginsburg, 2006; WIPO, 2016) delimitan con precisión la doctrina de los usos honrados y las excepciones para fines educativos y científicos.

\begin{table}[htbp]
\centering
\small
\renewcommand{\arraystretch}{1.25}
\textbf{Tabla 14.2}\\[0.15cm]
\textit{Matriz Comparativa entre la Regulación Ética Académica y el Régimen Jurídico de Propiedad Intelectual}\\[0.25cm]
\begin{tabularx}{\textwidth}{p{3.2cm}X X}
\toprule
\textbf{Criterio de Análisis} & \textbf{Dimensión Ética Académica} & \textbf{Régimen Jurídico (Derecho de Autor)} \\
\midrule
\textbf{Objeto Tutelado} & Ideas originales, hipótesis, metodologías, algoritmos, datos experimentales, texto e imágenes. & Expresión formal concreta y original de la obra; no tutela conceptos abstractos ni fórmulas en sí mismas. \\
\textbf{Instancia Juzgadora} & Docente titular, comité de ética de investigación, tribunal de grado y autoridades universitarias. & Poder Judicial, juzgados especializados en propiedad intelectual y órganos administrativos (e.g., Indecopi). \\
\textbf{Consecuencia Típica} & Desaprobación de la asignatura, reformulación integral del trabajo, anulación de grado o expulsión. & Indemnización de daños y perjuicios (ámbito civil); penas privativas de la libertad y multas (ámbito penal). \\
\textbf{Configuración Cruzada} & Una conducta puede cumplir la ley de autor (uso de una idea no protegida) y constituir falta ética grave por no citar. & Una conducta puede respetar la ética (citar la fuente) e infringir la ley por exceder los límites del uso honrado. \\
\bottomrule
\end{tabularx}
\vspace{0.15cm}
\raggedright \scriptsize \textbf{Nota.} Elaboración propia a partir del marco doctrinario de De Jesús G. (2016) y la legislación comparada.
\end{table}

\subsubsection{Derechos Morales y Patrimoniales en el Marco del Decreto Legislativo N.º 822 de Perú}

En el marco legal peruano, la regulación madre de la materia se encuentra consagrada en el \textbf{Decreto Legislativo N.º 822 (Ley sobre el Derecho de Autor)}, promulgado el 24 de abril de 1996 y armonizado con las directrices de la Organización Mundial de la Propiedad Intelectual (OMPI) y el Convenio de Berna. Esta normativa establece una distinción estructural de capital importancia para la praxis investigativa: la división entre \textbf{derechos morales} y \textbf{derechos patrimoniales} (Perú, 1996).

Los \textbf{derechos morales} (Artículos 21 al 29 del D.L. 822) representan el vínculo espiritual, inalienable y perpetuo que une al creador humano con su obra intelectual. La ley estipula categóricamente que estos derechos son \textit{perpetuos, inalienables, inembargables, irrenunciables e imprescriptibles}. Dentro de esta esfera destacan dos facultades angulares:
\begin{enumerate}
    \item \textbf{Derecho de Paternidad (Art. 22):} Faculta al autor a exigir en todo momento el reconocimiento de su condición de creador, exigiendo que su nombre o seudónimo sea consignado indefectiblemente cada vez que la obra sea reproducida, divulgada, traducida o adaptada. En el ámbito académico, este derecho sustenta la obligación indeclinable de citar las fuentes primarias.
    \item \textbf{Derecho de Integridad (Art. 23):} Otorga al autor la facultad jurídica de oponerse a cualquier deformación, mutilación, modificación o alteración de su obra que perjudique su reputación científica o menoscabe el sentido original de su discurso.
    \item \textbf{Derechos de Divulgación, Modificación y Retiro (Arts. 24--26):} Determinan la facultad exclusiva de decidir si la obra ha de ser publicada, modificarla antes o después de divulgada, o retirarla de la circulación comercial previa indemnización de perjuicios.
\end{enumerate}

Por su parte, los \textbf{derechos patrimoniales} (Artículos 31 al 40 del D.L. 822) conciernen a la dimensión económica y de explotación comercial de la creación. A diferencia de los derechos morales, los derechos patrimoniales son \textit{transferibles, embargables, temporalmente delimitados y renunciables}. Comprenden la potestad exclusiva de autorizar o prohibir la reproducción total o parcial de la obra por cualquier medio técnico (Art. 32), su comunicación pública (Art. 33), su distribución mediante venta o alquiler (Art. 34) y su traducción, adaptación o arreglo instrumental (Art. 35).

\begin{table}[htbp]
\centering
\small
\renewcommand{\arraystretch}{1.25}
\textbf{Tabla 14.3}\\[0.15cm]
\textit{Estructura Comparativa de Derechos Morales y Patrimoniales bajo el Decreto Legislativo N.º 822}\\[0.25cm]
\begin{tabularx}{\textwidth}{p{3.0cm}X X}
\toprule
\textbf{Eje de Comparación} & \textbf{Derechos Morales (Arts. 21--29 D.L. 822)} & \textbf{Derechos Patrimoniales (Arts. 31--40 D.L. 822)} \\
\midrule
\textbf{Bien Jurídico Tutelado} & Vínculo personalísimo, ético e intelectual del autor con la creación de su intelecto. & Monopolio de explotación económica y aprovechamiento financiero derivado de la obra. \\
\textbf{Naturaleza Jurídica} & Perpetuos, inalienables, inembargables, irrenunciables e imprescriptibles (Art. 21). & Transmisibles por causa de muerte, cesibles inter vivos, temporales y embargables (Art. 31). \\
\textbf{Facultades Principales} & Paternidad (Art. 22), Integridad (Art. 23), Divulgación (Art. 24), Modificación y Retiro (Art. 25). & Reproducción (Art. 32), Comunicación pública (Art. 33), Distribución (Art. 34), Traducción y Adaptación (Art. 35). \\
\textbf{Implicancia en la Universidad} & Fundamenta la obligatoriedad imperativa de la cita formal y prohíbe las autorías espurias o usurpadas. & Explica por qué el estudiante o investigador no puede reproducir capítulos enteros sin licencia de la editorial. \\
\bottomrule
\end{tabularx}
\vspace{0.15cm}
\raggedright \scriptsize \textbf{Nota.} Elaboración propia a partir del texto del Decreto Legislativo N.º 822 (Perú, 1996).
\end{table}

Esta distinción aclara un error habitual en el entorno de posgrado: cuando un investigador firma un contrato de cesión de derechos patrimoniales con una casa editorial internacional (como Springer, Elsevier o IEEE), cede el derecho exclusivo de comercialización e impresión masiva, pero jamás pierde su condición de autor ni su derecho moral de paternidad. En correlato, ningún usuario de ese artículo queda eximido del deber moral y legal de citar formalmente al autor original, aun cuando la descarga del texto se efectúe a través de una base de datos institucional licenciada por la universidad.

\subsubsection{Doctrina de Usos Honrados, Derecho de Cita y Licenciamiento Abierto}

El ordenamiento jurídico reconoce que el monopolio del titular del derecho de autor no puede ser absoluto, pues entraría en colisión con el derecho fundamental a la educación, la investigación científica y la libertad de información. Por esta razón, el Decreto Legislativo N.º 822 establece en su Título IV un régimen exhaustivo de \textbf{límites y excepciones al derecho patrimonial}.

El eje articulador de estas excepciones en el entorno académico es el \textbf{derecho de cita}, estipulado en el Artículo 44, inciso a:
\begin{quote}
\small \textit{\textsl{\guillemotleft Es lícita la reproducción por medio de la reprografía, para la enseñanza o la realización de exámenes en instituciones educativas, siempre que no haya fines de lucro y en la medida justificada por el objetivo perseguido, de artículos, breves extractos de obras o del conjunto de una obra plástica o fotográfica aislada, lícitamente publicadas, a condición de que tal utilización se haga conforme a los usos honrados y que la misma no sea objeto de venta u otra transacción a título oneroso, ni tenga directa o indirectamente fines de lucro\guillemotright}} (Perú, 1996, Art. 44).
\end{quote}

Asimismo, en el ámbito internacional, el Convenio de Berna para la Protección de las Obras Literarias y Artísticas (Art. 10, numeral 1) y la doctrina anglosajona del \textit{Fair Use} (Título 17, Sección 107 del Código de los Estados Unidos) configuran criterios concurrentes que delimitan cuándo la utilización de una obra ajena es jurídicamente admisible sin recabar autorización contractual expresa del titular:
\begin{itemize}
    \item \textbf{Propósito y Carácter del Uso:} El uso debe perseguir una finalidad de crítica académica, comentario ilustrativo, docencia sin ánimo de lucro o investigación científica transformadora.
    \item \textbf{Naturaleza de la Obra Protegida:} El uso de obras factuales, técnicas o científicas goza de mayor amplitud que el de creaciones artísticas de ficción pura.
    \item \textbf{Cantidad y Sustancialidad del Fragmento Reutilizado:} La porción extraída debe ser cuantitativa y cualitativamente proporcional al objetivo del análisis. Citar una o dos páginas de un volumen monográfico de trescientas páginas para criticar su metodología es lícito; reproducir el 40\% del texto vulnera el uso honrado.
    \item \textbf{Efecto sobre el Mercado Potencial de la Obra Original:} La reutilización no puede competir comercialmente con la obra original ni constituir un sustituto de mercado que desincentive su adquisición legítima.
\end{itemize}

\begin{figure}[htbp]
\centering
\begin{tikzpicture}[
    node distance=1.5cm and 2.0cm,
    every node/.style={font=\footnotesize\sffamily, align=center},
    layer/.style={rectangle, draw=unmsm_red, line width=0.8pt, fill=bg_marfil, rounded corners=3pt, text width=12cm, inner sep=8pt},
    sublayer/.style={rectangle, draw=unmsm_darkred, line width=0.6pt, fill=bg_arena, rounded corners=2pt, text width=10.5cm, inner sep=6pt},
    core/.style={rectangle, draw=unmsm_red, line width=1pt, fill=unmsm_red, text=white, rounded corners=2pt, text width=9.0cm, inner sep=6pt}
]
    \node [layer] (l1) {
        \textbf{CULTURA DE INTEGRIDAD CIENTÍFICA INSTITUCIONAL (UNMSM)}\\
        \scriptsize Marco Axiológico ICAI: Honestidad $\cdot$ Confianza $\cdot$ Justicia $\cdot$ Respeto $\cdot$ Responsabilidad $\cdot$ Valentía
    };
    \node [sublayer, below=0.4cm of l1] (l2) {
        \textbf{RÉGIMEN ÉTICO DISCIPLINAR UNIVERSITARIO}\\
        \scriptsize Evaluación de Procesos $\cdot$ Trazabilidad de Bitácoras $\cdot$ Rúbricas Analíticas $\cdot$ Criterios ICMJE / CRediT
    };
    \node [sublayer, below=0.4cm of l2] (l3) {
        \textbf{MARCO NORMATIVO LEGAL (D.L. 822 / TRATADOS OMPI)}\\
        \scriptsize Derechos Morales (Paternidad e Integridad) $\cdot$ Derechos Patrimoniales (Reproducción y Distribución)
    };
    \node [core, below=0.4cm of l3] (l4) {
        \textbf{DOCTRINA DE USOS HONRADOS Y EXCEPCIONES}\\
        \scriptsize Derecho de Cita (Art. 44) $\cdot$ Licencias Creative Commons $\cdot$ Acceso Abierto
    };
    
    \draw [->, thick, draw=unmsm_red] (l1) -- (l2);
    \draw [->, thick, draw=unmsm_red] (l2) -- (l3);
    \draw [->, thick, draw=unmsm_red] (l3) -- (l4);
\end{tikzpicture}
\caption*{\textbf{Figura 14.1.} \textit{Jerarquía de Protección Jurídica y Ética de la Producción Intelectual Universitaria.} \scriptsize \textbf{Fuente:} Elaboración propia basada en D.L. 822 y directrices ICAI (2021).}
\end{figure}

Un fenómeno colateral de gran relevancia en la docencia moderna es la proliferación de obras amparadas bajo esquemas de \textbf{Acceso Abierto} (\textit{Open Access}) y licencias estandarizadas como \textbf{Creative Commons (CC)}. Es un error recurrente entre estudiantes asumir que un documento de acceso abierto constituye material en el dominio público susceptible de cualquier explotación arbitraria. Como advierte la American Psychological Association (APA, 2022), el acceso abierto garantiza la gratuidad de la descarga y lectura, pero cada recurso permanece rigurosamente gobernado por los términos de su licencia:
\begin{itemize}
    \item \textbf{CC BY (Atribución):} Permite redistribuir, modificar y construir a partir de la obra, incluso comercialmente, exigiendo siempre dar crédito formal al creador.
    \item \textbf{CC BY-SA (Compartir Igual):} Exige que cualquier obra derivada sea licenciada bajo términos idénticos a los del original.
    \item \textbf{CC BY-NC (No Comercial):} Prohíbe terminantemente el aprovechamiento económico directo o indirecto del contenido.
    \item \textbf{CC BY-ND (Sin Derivadas):} Permite la redistribución únicamente en su forma íntegra e inalterada, prohibiendo adaptaciones o traducciones no autorizadas.
\end{itemize}

En concordancia con las normas APA (2022), cuando un tesista o investigador desea reproducir de forma exacta o adaptar un elemento visual de alta complejidad (un diagrama de flujo extenso, una escala psicométrica validada o una infografía arquitectónica), la mera citación parentética resulta insuficiente: es mandatorio consignar en la nota al pie de la figura que el material ha sido reproducido con el permiso formal del titular de los derechos o al amparo de una licencia CC específica, detallando el enlace a la licencia correspondiente.

La tensión entre la privatización de los activos intelectuales y la democratización del conocimiento en la era digital fue examinada críticamente por Lawrence Lessig (2001, 2004) en \textit{The Future of Ideas} y \textit{Free Culture}, así como por Siva Vaidhyanathan (2001) al alertar sobre el avance desmedido de los monopolios editoriales. En respuesta a estas restricciones comerciales, el movimiento de Acceso Abierto (\textit{Open Access}), teorizado por Peter Suber (2012), postula que la literatura científica debe ser un bien común accesible sin tarifas de suscripción, siempre que se garantice la atribución debida. En el ámbito nacional, guías institucionales del Indecopi (2018) y análisis sobre integridad académica en el Perú (Saldaña, 2021) subrayan que el aprovechamiento de repositorios abiertos fortalece la soberanía científica universitaria en estricto apego a la Ley Universitaria (Perú, 2014) y al D.L. 822.

\subsubsection{La Reproducibilidad Computacional como Imperativo Deontológico}

En el marco específico de la Facultad de Ciencias Matemáticas y la computación científica, el debate sobre la integridad académica trasciende la superficie del texto para alcanzar los artefactos algorítmicos. La formulación de un modelo predictivo, la resolución numérica de una ecuación diferencial estocástica o el entrenamiento de un clasificador neuronal no constituyen procesos puramente formales: involucran código fuente, conjuntos de datos y entornos de ejecución.

Como argumentan Palacios Carrasco et al. (2026), documentar con meticulosidad quirúrgica el pipeline experimental completo es una manifestación directa del principio de honestidad científica. La integridad computacional impone la obligatoriedad de garantizar la \textbf{reproducibilidad}: cualquier par académico con acceso a los datos y al entorno de compilación debe estar en capacidad de ejecutar el código y obtener exactamente las mismas métricas de rendimiento y resultados numéricos informados en el texto.

\begin{table}[htbp]
\centering
\small
\renewcommand{\arraystretch}{1.25}
\textbf{Tabla 14.4}\\[0.15cm]
\textit{Protocolo de Integridad y Reproducibilidad para Artefactos Computacionales en Ciencias Exactas}\\[0.25cm]
\begin{tabularx}{\textwidth}{p{2.8cm}X p{5.5cm}}
\toprule
\textbf{Artefacto de I+D} & \textbf{Requisito Deontológico de Transparencia y Atribución} & \textbf{Riesgo Técnico y Epistémico si se Omite} \\
\midrule
\textbf{Código Fuente} & Declarar en el encabezado de cada script o en el archivo \texttt{README.md} el autor, la procedencia del repositorio original, la fecha de acceso y la licencia de software (MIT, Apache 2.0, GPLv3). & Apropiación indebida de soluciones algorítmicas de terceros; vulneración de licencias de software libre. \\
\textbf{Conjuntos de Datos (\textit{Datasets})} & Citar la fuente primaria del repositorio de datos (Kaggle, Zenodo, UCI), la versión del archivo, la fecha de extracción y la descripción formal de los scripts de filtrado o imputación de valores faltantes. & Manipulación no documentada de datos; resultados espurios irreproducibles; sesgos de selección inadvertidos. \\
\textbf{Modelos Numéricos y Redes Neuronales} & Registrar de forma determinista las semillas aleatorias (\textit{random seeds}), hiperparámetros de convergencia, tasas de aprendizaje y librerías utilizadas junto con sus versiones precisas en un archivo \texttt{requirements.txt}. & Fluctuación estocástica de resultados; incapacidad de verificar empíricamente la estabilidad de los algoritmos. \\
\textbf{Figuras y Gráficos Sintéticos} & Elaborar visualizaciones vectoriales nativas mediante librerías reproducibles (Matplotlib, TikZ) o incluir la fuente formal, el permiso de reproducción y la licencia de reutilización en la nota al pie. & Plagio visual por copia de diagramas; infracción de derechos patrimoniales editoriales sobre esquemas. \\
\bottomrule
\end{tabularx}
\vspace{0.15cm}
\raggedright \scriptsize \textbf{Nota.} Elaboración propia a partir de las directrices metodológicas de Grupo 9 (Palacios Carrasco et al., 2026).
\end{table}

\subsubsection{Caso Práctico Ilustrativo 11.1: Reutilización de Figuras Complejas en Informes Técnicos}

\subsubsection*{Contexto:}
Un estudiante de la carrera de Computación Científica se encuentra finalizando su monografía de graduación sobre optimización combinatoria aplicada a redes de distribución logística. En el Capítulo 2 (Marco Teórico), el estudiante identifica en un artículo publicado en una prestigiosa revista indexada de la editorial Springer un esquema vectorial de alta densidad que resume la arquitectura matemática del algoritmo genético que pretende implementar. El artículo científico original fue publicado bajo un régimen comercial tradicional de reserva absoluta de derechos (\textit{All rights reserved}).

\subsubsection*{Análisis e Implementación:}
El estudiante se enfrenta a una encrucijada técnica y jurídica. Una lectura superficial de la doctrina del \textit{derecho de cita} podría inducirlo a creer que recortar la figura como captura de pantalla, pegarla en su documento en formato \LaTeX\ y añadir la leyenda \textsl{\guillemotleft Figura 2.3. Tomado de Hernández et al. (2024; caso de referencia hipotético)\guillemotright} es una práctica admisible. Sin embargo, en la dogmática jurídica y editorial estricta, las figuras complejas e ilustraciones gozan de protección autónoma como obras plásticas o gráficas aisladas; reproducir la totalidad de una figura sin licencia constituye una infracción del derecho patrimonial de reproducción de la editorial titular.

Para proceder con estricta rigurosidad, el estudiante adopta la siguiente alternativa metodológica legítima:
\begin{enumerate}
    \item \textbf{Opción A (Solicitud de Licencia Expresa):} Tramita a través del sistema automatizado \textit{RightsLink} de la editorial una licencia formal de reutilización no comercial para tesis universitaria (la cual, en la gran mayoría de las casas académicas, se concede de manera gratuita e instantánea en formato digital). Con la licencia aprobada, incorpora la figura añadiendo en la nota al pie: \textsl{\guillemotleft Reproducido con autorización de Springer Nature de Hernández et al. (2024), Licencia N.º 58210394821\guillemotright}.
    \item \textbf{Opción B (Síntesis y Elaboración Propia en TikZ):} En lugar de reproducir servilmente la imagen externa, el estudiante re-diseña conceptualmente el esquema empleando el paquete vectorial \texttt{tikz} en \LaTeX. El nuevo gráfico refleja su propia interpretación de las fases del algoritmo, adapta la nomenclatura matemática a las variables de su tesis y añade en el rótulo inferior: \textsl{\guillemotleft Figura 2.3. Diagrama de flujo del algoritmo genético modificado. Fuente: Elaboración propia a partir de los fundamentos teóricos propuestos por Hernández et al. (2024; caso de referencia hipotético)\guillemotright}.
\end{enumerate}

\subsubsection*{Interpretación y Resultados:}
Al optar por la reelaboración conceptual vectorizada (Opción B), el estudiante no solo disipa todo riesgo de infracción patrimonial de derechos de autor, sino que ejecuta un acto cognitivo superior: interioriza la mecánica del algoritmo, la rediseña conforme a las necesidades de su investigación y genera un recurso visual original que aporta valor genuino a la literatura de su facultad.

\subsection{Taxonomía Exhaustiva del Plagio y Patologías del Reconocimiento Autoral}

\subsubsection{Clasificación Morfológica del Plagio Académico}

El fenómeno del plagio en el contexto universitario contemporáneo se manifiesta a través de un espectro morfológico diverso que trasciende con creces el arquetipo rudimentario de la copia mecánica (\textit{copy-paste}). En la literatura especializada y en los reglamentos de integridad académica más avanzados, se identifican con claridad cinco modalidades estructurales:

\begin{enumerate}
    \item \textbf{Plagio Literal o Servil:} Consiste en la transposición exacta de un bloque discursivo ajeno sin emplear comillas tipográficas, sin formato de párrafo en bloque sangrado y sin consignar la cita parentética con autor, año y página. En el ámbito técnico, se manifiesta también como la copia íntegra de subrutinas de código fuente o de tablas de resultados experimentales.
    \item \textbf{Plagio Mosaico o Patchwork Plagiarism (Parafraseo Inadecuado):} Constituye una de las patologías más extendidas en la educación superior. Consiste en una reformulación superficial consistente en alterar el orden secuencial de un par de oraciones o sustituir mecánicamente ciertas palabras por sinónimos de uso corriente, conservando intacta la estructura argumentativa, la línea de deducción y el entramado lógico de la fuente matriz, sin brindar el crédito correspondiente. Cambiar un adjetivo o invertir la subordinación gramatical no convierte una idea ajena en un aporte original.
    \item \textbf{Autoplagio o Reciclaje Textual no Declarado:} De acuerdo con las pautas de la American Psychological Association (APA, 2022), el autoplagio ocurre cuando un estudiante o investigador presenta como si fuera inédito y original un trabajo, sección o conjunto de datos que ya fue sometido a evaluación académica, calificado en una asignatura previa o publicado formalmente en un acta de congreso. Esta práctica desvirtúa la finalidad evaluativa de las asignaturas universitarias (que exige certificar el desarrollo de competencias actuales) y contraviene la ética científica al inflar artificialmente la producción intelectual.
    \item \textbf{Plagio Estructural y Metodológico:} Se configura mediante la apropiación no declarada del diseño experimental, la arquitectura analítica o la secuencia de inferencias formales ideadas por otro investigador, aun cuando el infractor redacte el texto íntegramente con palabras distintas. En ciencias exactas, replicar la derivación paso a paso de un teorema sin reconocer la obra pionera que trazó dicho camino constituye una omisión inaceptable.
    \item \textbf{Plagio en Código Fuente y Datos Computacionales:} Específico de las disciplinas informáticas, comprende la práctica de descargar scripts de repositorios ajenos o de compañeros de curso, modificando cosméticamente los nombres de las variables, reordenando las funciones o alterando la indentación para eludir herramientas de análisis estático como \textit{MOSS} (\textit{Measure of Software Similarity}) o \textit{JPlag}, pretendiendo que la lógica del programa fue concebida por el infractor.
\end{enumerate}

\begin{table}[htbp]
\centering
\small
\renewcommand{\arraystretch}{1.25}
\textbf{Tabla 14.5}\\[0.15cm]
\textit{Taxonomía Multidimensional de Tipologías de Plagio en la Educación Superior}\\[0.25cm]
\begin{tabularx}{\textwidth}{p{2.8cm}X p{5.5cm}}
\toprule
\textbf{Modalidad de Plagio} & \textbf{Mecanismo de Ejecución y Alteración Discursiva} & \textbf{Señal de Alerta Técnica y Detección Forense} \\
\midrule
\textbf{Plagio Literal} & Copia textual exacta de pasajes sin comillas ni identificación de autor, año y página. & Discontinuidad abrupta de estilo, vocabulario y sintaxis entre párrafos contiguos del informe. \\
\textbf{Plagio Mosaico} & Reemplazo cosmético de sinónimos manteniendo la sintaxis y secuencia lógica original sin cita. & Coincidencia cuasi-exacta de la cadena temática y ejemplos con una fuente no declarada en el texto. \\
\textbf{Autoplagio} & Reutilización de informes calificados en ciclos anteriores presentándolos como nuevos. & Presencia de formulaciones o datos de asignaturas pasadas sin consultar al docente ni citar el trabajo base. \\
\textbf{Plagio Estructural} & Apropiación del diseño experimental, modelo abstracto o prueba de teoremas sin autoría. & Metodología idéntica a una tesis reciente cambiando únicamente el dominio muestral superficial. \\
\textbf{Plagio de Código} & Renombramiento de identificadores y cambio cosmético de bucles para eludir análisis estático. & Grafo de flujo de control y llamadas a funciones idéntico a una solución pública en GitHub. \\
\bottomrule
\end{tabularx}
\vspace{0.15cm}
\raggedright \scriptsize \textbf{Nota.} Elaboración propia basada en Palacios Carrasco et al. (2026), APA (2022) e ICAI (2021).
\end{table}

\subsubsection{El Parafraseo Legítimo frente al Mosaico: Protocolo de Transposición Cognitiva}

El antídoto metodológico contra el plagio mosaico reside en el dominio técnico del \textbf{parafraseo legítimo}. Parafrasear no es buscar sinónimos en un procesador de textos ni mezclar retazos de dos fuentes distintas; parafrasear es un proceso intelectual de alto orden que exige \textit{comprender, interiorizar, sintetizar y reconstruir} la idea ajena integrándola armónicamente dentro del propio discurso argumentativo, sin omitir bajo ninguna circunstancia la cita formal de la fuente original.

La teorización pedagógica del plagio mosaico y el fenómeno del \textit{patchwriting} fue desarrollada pioneramente por Rebecca Moore Howard (1995, 1999), quien demostró que el ensamblaje de fragmentos textuales suele constituir una fase transicional del aprendizaje cognitivo en la apropiación de discursos científicos complejos, y no necesariamente una conducta fraudulenta dolosa. En concordancia con ello, Diane Pecorari (2003, 2008) y Swales y Feak (2012) destacan que la adquisición de competencias de escritura académica avanzada exige una instrucción deliberada en síntesis retórica y citación transparente. Asimismo, Richard Posner (2007) en \textit{The Little Book of Plagiarism} delimita con precisión la influencia legítima del plagio como engaño al lector, mientras que Miguel Roig (2001, 2015) ofrece directrices internacionales para erradicar el autoplagio y la duplicación no declarada de textos.

El protocolo formativo propuesto por el Grupo 9 (Palacios Carrasco et al., 2026) articula cinco fases operativas:
\begin{enumerate}
    \item \textbf{Lectura Comprensiva Profunda:} Leer el pasaje de la fuente primaria hasta asimilar cabalmente su núcleo conceptual y los supuestos bajo los cuales opera.
    \item \textbf{Ocultamiento del Texto Matriz:} Apartar la vista de la fuente original para impedir que su estructura sintáctica condicione la redacción propia.
    \item \textbf{Redacción en Voz Autoral Propia:} Redactar la idea empleando la propia terminología técnica y el encuadre lógico del trabajo que se elabora.
    \item \textbf{Cotejo de Divergencia Estructural:} Contrastar la versión producida con el original para certificar que no persisten secuencias idénticas de más de tres o cuatro palabras clave comunes y que no se ha desvirtuado el sentido original.
    \item \textbf{Inserción de la Atribución Formal:} Incorporar la cita parentética completa (autor, año) o narrativa, especificando el pasaje específico según las directrices de la APA 7.ª edición.
\end{enumerate}

Para evidenciar la diferencia cualitativa entre la copia servil, el mosaico y el parafraseo genuino, analicemos el siguiente texto formal redactado para fines didácticos en este capítulo:

\begin{quote}
\small \textbf{Texto Fuente Original:}\\
\textsl{\guillemotleft La correlación matemática observada entre dos variables empíricas no demuestra una relación causal unívoca, puesto que ambas magnitudes pueden depender de un tercer factor subyacente no considerado en el modelo original\guillemotright} (Palacios Carrasco et al., 2026, p. 7).
\end{quote}

A continuación se presentan las diversas respuestas observadas en la práctica evaluativa:
\begin{itemize}
    \item \textbf{Versión A (Plagio Literal):} \textsl{\guillemotleft La correlación matemática observada entre dos variables empíricas no demuestra una relación causal unívoca, puesto que ambas magnitudes pueden depender de un tercer factor subyacente no considerado en el modelo original\guillemotright}.\\
    \textbf{Diagnóstico Forense:} Plagio puro y simple. Transposición mecánica sin comillas ni referencia bibliográfica.
    \item \textbf{Versión B (Plagio Mosaico):} \textsl{\guillemotleft El vínculo matemático registrado entre dos magnitudes empíricas no prueba una conexión de causa directa, debido a que las dos variables pueden estar supeditadas a un tercer elemento oculto no incluido en el diseño\guillemotright}.\\
    \textbf{Diagnóstico Forense:} Plagio mosaico. Se sustituyeron sinónimos cosméticos (\textsl{correlación $\rightarrow$ vínculo, magnitudes $\rightarrow$ variables, factor subyacente $\rightarrow$ elemento oculto}), pero la sintaxis, el orden de las cláusulas y el razonamiento exacto pertenecen a la fuente primaria, la cual ha sido omitida intencionalmente.
    \item \textbf{Versión C (Parafraseo Legítimo con Citación Rigurosa):} \textsl{\guillemotleft En el diseño de modelos estadísticos, la simple covariación cuantitativa entre variables no autoriza a inferir causalidad determinista, dado que pueden actuar variables confusoras invisibles que expliquen la variación conjunta de los datos (Palacios Carrasco et al., 2026)\guillemotright}.\\
    \textbf{Diagnóstico Forense:} Parafraseo legítimo ejemplar. El estudiante reelabora el principio con su propio léxico disciplinar (\textsl{covariación, variables confusoras}), demuestra dominio del concepto y reconoce de forma impecable la fuente teórica de la que procede la premisa.
\end{itemize}

\subsubsection{Instrumentos Operativos de Prevención: La Bitácora de Lectura Científica}

La prevención del plagio involuntario (aquel derivado de la desorganización documental y las prisas previas a la fecha límite) exige la adopción sistemática de herramientas de gestión del conocimiento. Como señalan Palacios Carrasco et al. (2026), los autores proponen la implementación obligatoria de una \textbf{Bitácora de Lectura Científica} antes de escribir la primera línea del borrador final.

La bitácora constituye un repositorio estructurado donde el investigador registra de forma inmutable los metadatos de cada fuente consultada, transcribiendo los pasajes clave estrictamente entrecomillados y elaborando de inmediato una paráfrasis preliminar junto con su valoración crítica. Al momento de redactar el manuscrito formal, cada enunciado factual puede rastrearse con exactitud matemática hasta su origen documental.

\begin{table}[htbp]
\centering
\small
\renewcommand{\arraystretch}{1.25}
\textbf{Tabla 14.6}\\[0.15cm]
\textit{Modelo Operativo de Bitácora de Lectura Científica para Ciencias Exactas}\\[0.25cm]
\begin{tabularx}{\textwidth}{p{3.2cm}p{1.8cm}p{2.2cm}X p{2.8cm}}
\toprule
\textbf{Fuente Primaria (APA 7 / DOI)} & \textbf{Ubicación} & \textbf{Tipo de Registro} & \textbf{Contenido Extraído o Paráfrasis Preliminar} & \textbf{Destino y Uso Previsto} \\
\midrule
De Jesús G. (2016). DOI: 10.5281/... & p. 48, párr. 3 & Cita textual exacta & \textsl{\guillemotleft El 80\% de los trabajos de grado revisados mostraron omisiones de fuentes primarias\guillemotright} & Marco teórico; justificación del problema. \\
Porto Castro (2024). Editorial Redipe. & pp. 228--230 & Paráfrasis técnica & Estudio cuantitativo sobre 310 universitarios donde se analiza la correlación entre procrastinación y plagio. & Discusión de antecedentes empíricos. \\
Miao \& Holmes (2023). UNESCO. & p. 15, sec. 2.1 & Resumen sintético & Las herramientas de IA deben subordinarse a la agencia humana; prohibición de autorías de LLMs. & Marco de gobernanza institucional. \\
\bottomrule
\end{tabularx}
\vspace{0.15cm}
\raggedright \scriptsize \textbf{Nota.} Modelo propuesto por Palacios Carrasco et al. (2026).
\end{table}

La articulación de esta bitácora permite consolidar el denominado \textbf{Ciclo de Trabajo Íntegro}, estructurado en siete etapas secuenciales que generan evidencia auditable a lo largo de todo el semestre académico:

\begin{figure}[htbp]
\centering
\begin{tikzpicture}[
    node distance=1.1cm and 0.5cm,
    font=\footnotesize\sffamily,
    box/.style={rectangle, draw=unmsm_red, line width=0.8pt, fill=bg_marfil, rounded corners=2pt, text width=2.2cm, text centered, minimum height=1.3cm},
    arrow/.style={-{Latex[length=2.5mm]}, thick, draw=unmsm_darkred}
]
    \node [box] (n1) {\textbf{01}\\[0.1cm]Lectura y Comprensión};
    \node [box, right=of n1] (n2) {\textbf{02}\\[0.1cm]Registro en Bitácora};
    \node [box, right=of n2] (n3) {\textbf{03}\\[0.1cm]Paráfrasis Legítima};
    \node [box, right=of n3] (n4) {\textbf{04}\\[0.1cm]Citación APA 7};
    \node [box, below=1.0cm of n4] (n5) {\textbf{05}\\[0.1cm]Verificación de Datos};
    \node [box, left=of n5] (n6) {\textbf{06}\\[0.1cm]Declaración de IA};
    \node [box, left=of n6] (n7) {\textbf{07}\\[0.1cm]Defensa Oral};
    
    \draw [arrow] (n1) -- (n2);
    \draw [arrow] (n2) -- (n3);
    \draw [arrow] (n3) -- (n4);
    \draw [arrow] (n4) -- (n5);
    \draw [arrow] (n5) -- (n6);
    \draw [arrow] (n6) -- (n7);
\end{tikzpicture}
\caption*{\textbf{Figura 14.2.} \textit{El Ciclo de Trabajo Íntegro en Siete Fases Trazables.} \scriptsize \textbf{Fuente:} Adaptado de Palacios Carrasco et al. (2026).}
\end{figure}

\subsubsection{Patologías de la Coautoría y Fraude Autorial: Criterios ICMJE/APA}

Junto al plagio discursivo, una de las mayores vulneraciones a la justicia y honestidad académica consiste en la alteración fraudulenta de las autorías en investigaciones científicas y proyectos universitarios. Como advirtió Robert K. Merton (1973) en su sociología de la ciencia, el reconocimiento autoral legítimo y la prioridad del descubrimiento constituyen el núcleo motivacional del avance del saber; cuando la autoría se desnaturaliza bajo la presión del \textit{publish or perish}, se distorsiona la confianza de la comunidad epistémica (Macfarlane, 2017). Para salvaguardar la transparencia en la investigación, el \textit{Committee on Publication Ethics} (COPE, 2019) y las investigaciones empíricas sobre deshonestidad académica de Donald McCabe (2001, 2005; McCabe et al., 2012), Tracey Bretag (2013, 2016), Wendy Sutherland-Smith (2008) y Tricia Bertram Gallant (2008) enfatizan que el reconocimiento debe ajustarse estrictamente a la contribución intelectual real. De Jesús G. (2016) y Palacios Carrasco et al. (2026) tipifican cuatro patologías recurrentes:

\begin{enumerate}
    \item \textbf{Autoría por Regalo o Cortesía (\textit{Gift / Courtesy Authorship}):} Concesión graciosa de la coautoría a un colega, estudiante o familiar que no participó en la investigación, motivada por lazos de amistad, parentesco o benevolencia mal entendida.
    \item \textbf{Autoría por Intercambio o Negociación (\textit{Mutual Support / Quid Pro Quo}):} Acuerdos colusorios mediante los cuales dos o más investigadores se comprometen a incluir mutuamente sus firmas en sus respectivos manuscritos sin haber contribuido intelectualmente, con el exclusivo fin de duplicar artificialmente sus índices métricos curriculares.
    \item \textbf{Autoría Honoraria o por Estatus (\textit{Guest Authorship}):} Inclusión unilateral de una personalidad de elevado prestigio académico, un director de instituto o un jefe de departamento para explotar su nombre y elevar las probabilidades de aprobación en revistas de alto impacto o tribunales de tesis.
    \item \textbf{Autoría Fantasma (\textit{Ghost Authorship}):} Omisión deliberada del crédito autoral de quien efectivamente diseñó, analizó o redactó la investigación (como tesistas, becarios subordinados o redactores técnicos remunerados), privándolos de su derecho moral de paternidad. En la actualidad, esta figura abarca también la utilización encubierta y no declarada de modelos fundacionales de IA generativa para redactar sustantivamente el texto.
\end{enumerate}

Para erradicar estas distorsiones, la comunidad científica internacional acoge los cuatro criterios concurrentes fijados por el \textit{International Committee of Medical Journal Editors} (ICMJE, 2023), asumidos plenamente por la APA (2022) y el D.L. 822 de Perú:
\begin{enumerate}
    \item Contribución sustancial e insustituible a la concepción, diseño del estudio o a la adquisición, procesamiento, análisis e interpretación matemática de los datos.
    \item Redacción del borrador original del artículo o revisión crítica profunda de su contenido intelectual.
    \item Aprobación formal y explícita de la versión final del manuscrito que será sometida a evaluación o publicación.
    \item Capacidad de respuesta y rendición de cuentas (\textit{accountability}): compromiso de responder solidariamente por la precisión, validez e integridad de todas las partes de la obra.
\end{enumerate}

La ley y la doctrina enfatizan que \textit{la mera provisión de fondos financieros, la cesión de instalaciones de laboratorio, el suministro de datos brutos sin análisis o la simple corrección ortotipográfica no confieren mérito autoral alguno}; quienes presten este tipo de apoyos colaterales deben figurar formalmente en la sección de \textbf{Agradecimientos} (\textit{Acknowledgments}).

\begin{table}[htbp]
\centering
\small
\renewcommand{\arraystretch}{1.25}
\textbf{Tabla 14.7}\\[0.15cm]
\textit{Matriz Diagnóstica de Autorías Ficticias y Mecanismos de Neutralización Institucional}\\[0.25cm]
\begin{tabularx}{\textwidth}{p{2.8cm}X p{5.5cm}}
\toprule
\textbf{Tipología Ficticia} & \textbf{Configuración Ética y Mecanismo de Distorsión} & \textbf{Instrumento Institucional de Neutralización} \\
\midrule
\textbf{Autoría por Regalo} & Inclusión de personas sin contribución intelectual por afecto, complicidad o paternalismo docente. & Exigencia de declaración jurada vinculada a la taxonomía de roles CRediT en cada entrega. \\
\textbf{Autoría por Intercambio} & Pacto colusorio entre pares para inflar currículos citándose o firmándose recíprocamente. & Auditoría cruzada de proyectos y verificación de competencias en sustentaciones individuales. \\
\textbf{Autoría Honoraria} & Adición de directores o científicos sénior sin aporte real para blindar el manuscrito. & Entrevistas orales obligatorias donde cada firmante debe explicar su aporte teórico directo. \\
\textbf{Autoría Fantasma} & Exclusión de tesistas o estudiantes que ejecutaron el trabajo; o uso encubierto de IA. & Trazabilidad de versiones mediante sistemas de control git y repositorio de borradores firmados. \\
\bottomrule
\end{tabularx}
\vspace{0.15cm}
\raggedright \scriptsize \textbf{Nota.} Adaptado de De Jesús G. (2016) y Palacios Carrasco et al. (2026).
\end{table}

\subsubsection{Operacionalización mediante la Taxonomía CRediT}

Para dotar de transparencia técnica a la asignación de responsabilidades, la comunidad internacional ha estandarizado la \textbf{Taxonomía CRediT} (\textit{Contributor Roles Taxonomy}, Brand et al., 2015), reconocida formalmente bajo la norma ANSI/NISO Z39.104-2022. Esta taxonomía desglosa la labor científica en catorce roles funcionales específicos:

\begin{itemize}
    \item \textbf{Conceptualización (\textit{Conceptualization}):} Formulación y desarrollo de las metas, preguntas y objetivos centrales de la investigación.
    \item \textbf{Curación de Datos (\textit{Data Curation}):} Actividades de gestión, depuración y mantenimiento de los datos brutos.
    \item \textbf{Análisis Formal (\textit{Formal Analysis}):} Aplicación de técnicas estadísticas, matemáticas y algorítmicas sobre los datos.
    \item \textbf{Adquisición de Fondos (\textit{Funding Acquisition}):} Obtención del financiamiento que hizo viable el estudio.
    \item \textbf{Investigación (\textit{Investigation}):} Ejecución de los experimentos, pruebas de laboratorio o trabajo de campo.
    \item \textbf{Metodología (\textit{Methodology}):} Desarrollo o diseño formal de los procedimientos experimentales y modelos.
    \item \textbf{Administración del Proyecto (\textit{Project Administration}):} Gestión y coordinación de las fases de ejecución del equipo.
    \item \textbf{Recursos (\textit{Resources}):} Provisión de materiales de cómputo, servidores o reactivos de laboratorio.
    \item \textbf{Software (\textit{Software}):} Programación, diseño algorítmico e implementación de librerías y componentes.
    \item \textbf{Supervisión (\textit{Supervision}):} Guía metodológica, mentoría y liderazgo del equipo investigador.
    \item \textbf{Validación (\textit{Validation}):} Verificación de la replicabilidad, estabilidad y precisión de los resultados.
    \item \textbf{Visualización (\textit{Visualization}):} Preparación y diseño de figuras, gráficos interactivos y tablas complejas.
    \item \textbf{Redacción del Borrador Original (\textit{Writing -- Original Draft}):} Redacción de la primera versión integral del texto.
    \item \textbf{Redacción, Revisión y Edición (\textit{Writing -- Review \& Editing}):} Revisión crítica, adición de aportes y edición sustantiva.
\end{itemize}

\subsubsection{Caso Práctico 11.2: Resolución de Conflictos de Autoría en Computación Científica}

\subsubsection*{Contexto:}
En el marco de la asignatura de Inteligencia Artificial Aplicada de la UNMSM, un equipo conformado por cuatro estudiantes (A, B, C y D) debe presentar como entregable final un proyecto de investigación que comprende: el desarrollo de un modelo predictivo basado en redes neuronales convolucionales, el código fuente ejecutable debidamente documentado en un repositorio institucional y un artículo científico redactado en formato IEEE / APA. 

Al culminar el semestre, se suscitan severas discrepancias internas respecto a la nómina de autores:
\begin{itemize}
    \item El estudiante \textbf{A} diseñó la arquitectura del modelo matemático y codificó el 100\% del software en Python.
    \item El estudiante \textbf{B} formuló el marco teórico formal, revisó la literatura previa y redactó la totalidad del artículo monográfico.
    \item El estudiante \textbf{C} recopiló el conjunto de datos de un repositorio público y colaboró en el etiquetado y limpieza preliminar.
    \item El estudiante \textbf{D} no participó en las reuniones de trabajo, no escribió código ni redactó el informe; no obstante, financió las impresiones del documento empastado y realizó una revisión ortográfica superficial la noche previa a la sustentación.
\end{itemize}
Los cuatro estudiantes demandan ante el docente titular figurar como coautores principales en estricta igualdad de condiciones.

\subsubsection*{Análisis e Implementación:}
El docente y el comité evaluador aplican los principios de justicia académica del ICAI (2021), los criterios vinculantes del ICMJE (2023) y el Artículo 21 del D.L. 822 de la República del Perú. Se procede a auditar los aportes mediante una matriz formal de la Taxonomía CRediT:

\begin{table}[htbp]
\centering
\small
\renewcommand{\arraystretch}{1.2}
\textbf{Tabla 14.8}\\[0.15cm]
\textit{Auditoría de Contribuciones CRediT y Dictamen Autorial del Caso Práctico 11.2}\\[0.25cm]
\begin{tabularx}{\textwidth}{p{2.0cm}p{4.0cm}X p{3.2cm}}
\toprule
\textbf{Integrante} & \textbf{Roles CRediT Verificados} & \textbf{Naturaleza del Aporte Intelectual} & \textbf{Dictamen Autorial} \\
\midrule
\textbf{Estudiante A} & Metodología, Software, Análisis Formal, Validación. & Creación del artefacto algorítmico y experimentación. & \textbf{Coautor Principal} (1.er lugar). \\
\textbf{Estudiante B} & Conceptualización, Redacción borrador, Revisión. & Marco teórico y articulación científica del paper. & \textbf{Coautor Principal} (2.º lugar). \\
\textbf{Estudiante C} & Curación de Datos, Recursos. & Adquisición y limpieza de datos brutos. & \textbf{Coautor Secundario} (condicionado a defensa). \\
\textbf{Estudiante D} & Ningún rol intelectual sustancial. & Apoyo logístico, financiero y revisión ortográfica leve. & \textbf{Exclusión de Autoría}; mención en Agradecimientos. \\
\bottomrule
\end{tabularx}
\vspace{0.15cm}
\raggedright \scriptsize \textbf{Nota.} Análisis forense basado en los criterios ICMJE (2023) y Palacios Carrasco et al. (2026).
\end{table}

\subsubsection*{Interpretación y Resultados:}
El dictamen excluye categóricamente al estudiante D de la condición de autor. El financiamiento de materiales fungibles (como papel o tóner) y la corrección ortotipográfica no constituyen aportes intelectivos sustanciales al diseño de la investigación. Incluirlo hubiera tipificado una \textbf{autoría de cortesía fraudulenta}. Al estudiante C se le concede la coautoría secundaria supeditada a que demuestre en una sustentación oral el dominio sobre el preprocesamiento de datos y la defensa de las hipótesis del estudio. Este dictamen restaura la justicia académica en el aula y previene la propagación de vicios en la carrera profesional de los futuros científicos.

\subsection{Inteligencia Artificial Generativa, Postplagio y Transformación Epistémica}

\subsubsection{La Disrupción de los Modelos Fundacionales en la Producción Textual y de Código}

La eclosión masiva de los modelos fundacionales basados en la arquitectura \textit{Transformer} (Vaswani et al., 2017) y el aprendizaje por refuerzo con retroalimentación humana (\textit{RLHF}, Ouyang et al., 2022) ha introducido una disrupción copernicana en los procesos de escritura científica y codificación de software. Por primera vez en la historia de la técnica, los agentes computacionales no operan como meros procesadores pasivos de caracteres o clasificadores deterministas, sino que poseen la capacidad probabilística de generar lenguaje natural con una sintaxis intachable, articular cadenas de deducción teórica y sintetizar librerías completas de código algorítmico a partir de instrucciones textuales en lenguaje natural (\textit{prompts}).

Esta capacidad desestabiliza las premisas sobre las cuales se estructuró históricamente la evaluación académica. Durante siglos, la capacidad de redactar un texto estructurado con fluidez retórica se consideró una prueba inequívoca de la existencia de una comprensión conceptual previa en la mente del redactor. Hoy, dicha inferencia ha colapsado: un texto puede ostentar una coherencia formal impecable, vocabulario especializado y un tono altamente persuasivo sin que exista un ápice de asimilación cognitiva en quien presionó la tecla de generación.

\subsubsection{Uso Formativo frente a Uso Sustitutivo: Fraude por Subcontratación Algorítmica}

Para salvaguardar el rigor del aprendizaje universitario, la doctrina contemporánea delimita con estricta precisión conceptual la frontera entre el \textbf{uso formativo} y el \textbf{uso sustitutivo} de la inteligencia artificial generativa (Palacios Carrasco et al., 2026; Miao \& Holmes, 2023).

El \textbf{uso formativo} concibe a la IA como un tutor dialógico y un andamiaje cognitivo (\textit{scaffolding}). En este esquema, el estudiante conserva el control epistemológico y la dirección de la tarea: recurre al modelo para solicitar explicaciones alternativas de un concepto matemático abstracto que no logró comprender en el texto guía, pedir observaciones críticas sobre la claridad argumentativa de un borrador redactado por él mismo, o explorar posibles casos de borde (\textit{edge cases}) para poner a prueba la robustez de una función programada en su laboratorio. En todos estos supuestos, el estudiante ejerce un juicio crítico activo, discrimina qué sugerencias incorporar y puede justificar matemáticamente cada decisión ante sus evaluadores.

Por el contrario, el \textbf{uso sustitutivo} configura una modalidad avanzada de fraude académico que podemos categorizar como \textbf{fraude por subcontratación algorítmica} (\textit{algorithmic contract cheating}). En esta práctica, el estudiante ingresa la consigna de la tarea o el problema matemático en la ventana del asistente generativo, copia la respuesta sintética sin mediar análisis alguno y la presenta como producto de su propio esfuerzo reflexivo. En este escenario, la actividad intelectiva ha sido vaciada por completo: el estudiante actúa como un mero intermediario mecánico entre el algoritmo y el profesor, resultando imposible certificar el desarrollo de las competencias académicas que la universidad está llamada a validar.

\begin{table}[htbp]
\centering
\small
\renewcommand{\arraystretch}{1.25}
\textbf{Tabla 14.9}\\[0.15cm]
\textit{Matriz Comparativa entre Uso Formativo Admisible y Uso Sustitutivo Fraudulento de la IA}\\[0.25cm]
\begin{tabularx}{\textwidth}{p{3.2cm}X X}
\toprule
\textbf{Eje de Comparación} & \textbf{Uso Formativo Admisible (Andamiaje Cognitivo)} & \textbf{Uso Sustitutivo Fraudulento (Subcontratación)} \\
\midrule
\textbf{Dirección del Proceso} & El estudiante formula las hipótesis, redacta los borradores y utiliza la IA para someter a crítica su trabajo. & La IA genera la tesis central, el desarrollo argumental y el código; el estudiante solo transcribe. \\
\textbf{Carga Intelectual} & La agencia cognitiva, la síntesis y la toma de decisiones finales residen plenamente en el autor humano. & Se evade el esfuerzo reflexivo y la internalización conceptual mediante la delegación al algoritmo. \\
\textbf{Trazabilidad y Control} & El estudiante conserva el registro de los borradores, los \textit{prompts} y declara con honestidad las herramientas. & Se oculta el uso de la máquina, presentando la salida computacional como producto mental genuino. \\
\textbf{Capacidad de Defensa} & El estudiante es capaz de responder con solvencia sobre cada ecuación, algoritmo y conclusión de la obra. & El estudiante es incapaz de explicar los fundamentos lógicos o depurar los fallos del código entregado. \\
\bottomrule
\end{tabularx}
\vspace{0.15cm}
\raggedright \scriptsize \textbf{Nota.} Elaboración propia basada en Palacios Carrasco et al. (2026) y directrices UNESCO (Miao \& Holmes, 2023).
\end{table}

\begin{table}[htbp]
\centering
\small
\renewcommand{\arraystretch}{1.25}
\textbf{Tabla 14.10}\\[0.15cm]
\textit{Guía Operativa de Admisibilidad de IA Generativa según Tipología de Tarea Académica}\\[0.25cm]
\begin{tabularx}{\textwidth}{p{3.5cm}p{3.2cm}X}
\toprule
\textbf{Tipología de Actividad} & \textbf{Criterio de Admisibilidad} & \textbf{Condición Pedagógica y Operativa Exigida} \\
\midrule
\textbf{Comprensión Teórica} & \textbf{Admisible} (formativo) & Contrastar la respuesta del modelo con la bibliografía obligatoria del curso y reformular en palabras propias. \\
\textbf{Revisión de Borrador Propio} & \textbf{Admisible} (formativo) & El estudiante debe justificar documentalmente qué sugerencias de la IA aceptó y cuáles rechazó por inexactas. \\
\textbf{Optimización Sintáctica de Código} & \textbf{Admisible con condición} & Declarar la procedencia del fragmento y estar en capacidad de explicar la complejidad temporal asintótica en pizarra. \\
\textbf{Traducción de Borrador Propio} & \textbf{Admisible si el sílabo lo permite} & Declarar la herramienta en el apéndice de metodología y revisar personalmente la fidelidad técnica del texto traducido. \\
\textbf{Generación Directa del Ensayo} & \textbf{Inadmisible} (fraude total) & Constituye subcontratación algorítmica; anula la evaluación y activa el procedimiento disciplinar. \\
\bottomrule
\end{tabularx}
\vspace{0.15cm}
\raggedright \scriptsize \textbf{Nota.} Propuesta de los autores a partir de la práctica docente en Computación Científica (UNMSM).
\end{table}

\subsubsection{La Teoría del Postplagio y la Reconfiguración de la Originalidad}

La consolidación irreversible de estas tecnologías en la sociedad del conocimiento ha llevado a teóricos de la educación superior a replantear el concepto mismo de originalidad. En una rigurosa revisión sistemática centrada en instituciones de educación superior, Corona Domínguez y Alonso Nuñez (2026) introducen en el debate hispanoamericano la noción de \textbf{postplagio} (\textit{postplagiarism}), originalmente formulada por Sarah Elaine Eaton (2021).

El postplagio postula que la academia ha ingresado en una época poshumanista de autoría híbrida donde la producción textual se realiza crecientemente en conjunción con agentes artificiales. En este nuevo horizonte, la obsesión punitiva tradicional orientada a comparar cadenas idénticas de texto para certificar la \textsl{\guillemotleft originalidad pura\guillemotright} resulta metodológicamente obsoleta e inoperante. El foco ético de la universidad del siglo XXI debe desplazarse desde el producto lingüístico estático hacia la \textbf{agencia cognitiva humana}: la autenticidad ya no se mide por la ausencia absoluta de asistencia algorítmica, sino por la capacidad demostrable del estudiante de ejercer curaduría crítica, verificar empíricamente los datos, aportar rigor conceptual y responsabilizarse éticamente por las consecuencias del conocimiento suscrito.

Asimismo, la Guía de la UNESCO para la IA generativa en educación e investigación (Miao \& Holmes, 2023) consagra como principio rector vinculante la \textbf{primacía de la agencia humana}. Los sistemas de IA no son sujetos de derecho ni poseen conciencia moral; por consiguiente, jamás pueden ser acreditados como autores o coautores de una publicación científica. Quien firma un trabajo asume una responsabilidad ética y legal exclusiva e indelegable; transferir la culpa de un error analítico o de una fuente apócrifa a la máquina (\textsl{\guillemotleft la IA me lo sugirió\guillemotright}) constituye una capitulación ética incompatible con la investidura académica universitaria.

\subsubsection{Patologías Epistémicas de los LLMs: Alucinaciones Algorítmicas}

Uno de los riesgos técnicos más perniciosos de los modelos fundacionales radica en su naturaleza probabilística: son sistemas diseñados para predecir el siguiente \textit{token} estadísticamente más verosímil de acuerdo con sus parámetros de entrenamiento, no motores deterministas de verificación de la verdad. Como demuestran Maynez et al. (2020) y Ji et al. (2023), los modelos de lenguaje priorizan la coherencia superficial del texto generado por sobre la fidelidad factual frente al corpus de referencia, provocando las denominadas \textbf{alucinaciones algorítmicas}: la invención persuasiva de hechos, datos estadísticos, demostraciones matemáticas falaces y referencias bibliográficas apócrifas. Frente a este riesgo, directrices éticas internacionales (COPE, 2023; European Commission, 2019, 2024; WIPO, 2023) prohíben taxativamente otorgar autoría o coautoría a herramientas de IA, reafirmando la primacía de la responsabilidad humana.

En el ámbito universitario, el fenómeno de la cita alucinada reviste extrema gravedad. Un modelo generativo puede estructurar una referencia bibliográfica dotándola de una apariencia de absoluta legitimidad científica: combina el nombre de un investigador respetado en el área, un título verosímil con terminología vanguardista, un nombre de revista indexada plausible, año de publicación reciente y hasta un prefijo numérico con el formato estándar de un DOI (\textit{Digital Object Identifier}). No obstante, al auditar dicho identificador en los servidores globales de \textit{Crossref} o \textit{DataCite}, se descubre que la publicación no existe en la realidad empírica.

Para neutralizar este vector de contaminación epistémica, Palacios Carrasco et al. (2026) formulan el \textbf{Protocolo de Verificación Heurística en Cinco Pasos}, de obligatoria observancia para todo estudiante antes de incorporar una fuente bibliográfica en sus trabajos:
\begin{enumerate}
    \item \textbf{Localización Primaria:} Buscar el título y DOI exacto en bases de datos científicas arbitradas (\textit{Scopus, Web of Science, IEEE Xplore, SciELO, Dialnet}).
    \item \textbf{Autenticación de Metadatos:} Corroborar que los nombres de los autores, el volumen, número y rango de páginas coinciden milimétricamente con el registro editorial formal.
    \item \textbf{Lectura Directa del Pasaje:} Acceder al documento original y leer personalmente el fragmento correspondiente para verificar que la fuente sostiene efectivamente la afirmación que se pretende respaldar.
    \item \textbf{Atribución a la Fuente Primaria:} Citar el artículo científico real, jamás a la herramienta de software que lo sugirió en su ventana de chat.
    \item \textbf{Registro de Trazabilidad:} Asentar la confirmación del DOI en la bitácora de lectura, adjuntando la declaración de uso de IA correspondiente.
\end{enumerate}

\subsubsection{Caso Práctico 11.3: Detección y Auditoría Forense de Alucinaciones Bibliográficas}

\subsubsection*{Contexto:}
Un estudiante de maestría presenta el primer borrador de su proyecto de tesis en ciencias exactas. En el apartado de antecedentes internacionales, incluye tres citas bibliográficas que sustentan la superioridad cuantitativa de un nuevo algoritmo de regularización estocástica. Las citas presentan la siguiente estructura de simulación docente: \textsl{\guillemotleft Martínez-Sánchez et al. (2025; ejemplo hipotético de alucinación sintética). High-dimensional stochastic regularization in convex optimization. Journal of Applied Mathematics and Computation, 412, 108--125. DOI: \url{10.1016/j.amc.2025.10892}\guillemotright}.

\subsubsection*{Análisis e Implementación:}
El docente supervisor realiza una auditoría forense rápida consultando la API pública de \textit{Crossref} mediante un script de validación en Python:

\begin{tcolorbox}[colback=bg_marfil, colframe=unmsm_gold, arc=2pt, boxrule=0.5pt, top=0.2cm, bottom=0.2cm, left=0.3cm, right=0.3cm]
\small\ttfamily
\begin{verbatim}
import urllib.request
import json

doi = "10.1016/j.amc.2025.10892"
url = f"https://api.crossref.org/works/{doi}"

try:
    headers = {'User-Agent': 'UNMSM_Audit/1.0'}
    req = urllib.request.Request(url, headers=headers)
    with urllib.request.urlopen(req) as response:
        data = json.loads(response.read().decode())
        print("Auténtico:", data['message']['title'][0])
except urllib.error.HTTPError as e:
    print(f"Error {e.code}: El DOI no existe en los registros de Crossref.")
\end{verbatim}
\end{tcolorbox}

La consulta arroja un código HTTP 404 (Not Found). Al indagar en la base de datos de la revista \textit{Applied Mathematics and Computation}, se constata que el volumen 412 corresponde al año 2022 y aborda temáticas distintas, y que ningún autor con ese apellido ha publicado sobre regularización estocástica en dicho número. Convocado a una sesión de revisión, el estudiante admite haber solicitado a un asistente de IA \textsl{\guillemotleft tres artículos recientes con DOI que demuestren que mi método es superior\guillemotright}, copiando el resultado en su bibliografía sin comprobar la existencia real de los documentos.

\subsubsection*{Interpretación y Resultados:}
El caso tipifica una gravísima falta de integridad científica: la incorporación de evidencia bibliográfica apócrifa derivada de una alucinación algorítmica. Aunque no existió una intención dolosa deliberada de falsificar datos manualmente, la negligencia temeraria en la verificación vulnera el valor de la responsabilidad científica. Como medida formativa, se anuló la entrega parcial, se obligó al estudiante a rehacer la revisión basándose exclusivamente en artículos físicos o digitales leídos y contrastados en biblioteca, y se le instruyó formalmente en el protocolo de declaración de IA.

\begin{figure}[htbp]
\centering
\begin{tikzpicture}[
    node distance=1.4cm and 1.0cm,
    font=\footnotesize\sffamily,
    startstop/.style={rectangle, draw=unmsm_red, line width=1pt, fill=bg_arena, rounded corners=3pt, text width=5.5cm, text centered, inner sep=6pt},
    decision/.style={diamond, draw=unmsm_darkred, line width=0.8pt, fill=bg_marfil, text width=3.8cm, text centered, inner sep=2pt, aspect=1.8},
    process/.style={rectangle, draw=unmsm_red, line width=0.8pt, fill=bg_marfil, rounded corners=2pt, text width=4.5cm, text centered, inner sep=6pt},
    arrow/.style={-{Latex[length=2.5mm]}, thick, draw=unmsm_darkred}
]
    \node [startstop] (start) {\textbf{INICIO:}\\Se contempla utilizar una herramienta de IA generativa en una tarea académica};
    \node [decision, below=0.8cm of start] (d1) {\textquestiondown La consigna del docente o el sílabo permite el uso de IA?};
    \node [process, right=1.2cm of d1] (no1) {\textbf{USO PROHIBIDO}\\[0.1cm]Resolver la tarea con esfuerzo propio sin herramientas de IA.};
    \node [decision, below=1.0cm of d1] (d2) {\textquestiondown El uso sustituye mi razonamiento o la autoría del texto/código?};
    \node [process, right=1.2cm of d2] (no2) {\textbf{FRAUDE ACADÉMICO}\\[0.1cm]Subcontratación algorítmica inaceptable. Reformular enfoque.};
    \node [process, below=1.0cm of d2] (yes) {\textbf{USO FORMATIVO ADMISIBLE}\\[0.1cm]Andamiaje cognitivo: lluvia de ideas, depuración de errores o crítica.};
    \node [startstop, below=0.8cm of yes] (decl) {\textbf{DECLARACIÓN TRANSPARENTE}\\[0.1cm]Adjuntar plantilla: herramienta, prompt, cambios realizados y verificación.};

    \draw [arrow] (start) -- (d1);
    \draw [arrow] (d1) -- node[above]{\textbf{NO}} (no1);
    \draw [arrow] (d1) -- node[right]{\textbf{SÍ}} (d2);
    \draw [arrow] (d2) -- node[above]{\textbf{SÍ}} (no2);
    \draw [arrow] (d2) -- node[right]{\textbf{NO}} (yes);
    \draw [arrow] (yes) -- (decl);
\end{tikzpicture}
\caption*{\textbf{Figura 14.3.} \textit{Árbol de Decisión Algorítmico para el Uso Ético de Herramientas de IA Generativa.} \scriptsize \textbf{Fuente:} Adaptado de Palacios Carrasco et al. (2026).}
\end{figure}

\begin{table}[htbp]
\centering
\small
\renewcommand{\arraystretch}{1.25}
\textbf{Tabla 14.11}\\[0.15cm]
\textit{Modelo Institucional de Declaración de Uso de Inteligencia Artificial Generativa}\\[0.25cm]
\begin{tabularx}{\textwidth}{p{3.8cm}X}
\toprule
\textbf{Dimensión de Transparencia} & \textbf{Declaración Detallada del Estudiante / Investigador} \\
\midrule
\textbf{1. Herramienta y Versión} & Claude 3.5 Sonnet / ChatGPT-4o (OpenAI), ejecutado mediante interfaz web oficial. \\
\textbf{2. Propósito y Prompt} & \textsl{\guillemotleft Identificar posibles contraargumentos a mi hipótesis sobre la convergencia del método de gradiente conjugado en espacios de Hilbert\guillemotright}. \\
\textbf{3. Salida Generada} & El modelo sugirió considerar el caso de operadores no acotados y matrices mal condicionadas. \\
\textbf{4. Juicio y Modificación} & Se aceptó analizar matrices mal condicionadas; se rechazó la sugerencia de operadores no acotados por exceder el dominio acotado fijado en el Teorema 3. \\
\textbf{5. Método de Verificación} & Se demostró analíticamente el lema de convergencia recurriendo al texto de Luenberger (1997) y se validó el código numérico en GNU Octave. \\
\bottomrule
\end{tabularx}
\vspace{0.15cm}
\raggedright \scriptsize \textbf{Nota.} Plantilla estandarizada de transparencia propuesta por Grupo 9 (Palacios Carrasco et al., 2026).
\end{table}

\subsection{Análisis Crítico de la Literatura Empírica y Cautela Metodológica}

\subsubsection{Auditoría Empírica en Trabajos de Grado: El Estudio de De Jesús G. (2016)}

Para comprender las raíces estructurales del fraude y las omisiones de citación en el contexto universitario latinoamericano, resulta imperativo examinar investigaciones empíricas de campo. Una de las obras de mayor relevancia metodológica es la conducida por De Jesús G. (2016), quien ejecutó una rigurosa auditoría documental sobre una muestra de 15 trabajos de grado aprobados durante el período 2010--2011 en la Facultad de Ciencias Económicas y Sociales (FACES) de la Universidad de Los Andes (Venezuela).

El diseño metodológico de la autora se basó en un muestreo no probabilístico e intencional, seleccionando 15 tesis representativas de entre un universo global de 60 monografías que ya contaban con la aprobación unánime de los jurados examinadores y se hallaban depositadas en el repositorio institucional. Los resultados revelaron un hallazgo alarmante: \textbf{12 de los 15 trabajos de grado (el 80\% de la muestra analizada) presentaron omisiones inaceptables de fuentes primarias y secundarias}, tipificando diversas formas de plagio literal, mosaico y alteración de citas.

\begin{figure}[htbp]
\centering
\begin{tikzpicture}[
    font=\footnotesize\sffamily,
    bar/.style={rectangle, rounded corners=2pt, minimum height=0.9cm}
]
    \node [anchor=west] at (0, 1.2) {\textbf{Tesis con Omisiones Graves de Fuentes (80\%) --- 12 trabajos}};
    \draw [fill=unmsm_red, draw=unmsm_darkred, line width=0.8pt] (0, 0.4) rectangle (9.6, 1.0);
    \node at (4.8, 0.7) [text=white, font=\small\sffamily\bfseries] {80\% de la muestra auditada};
    
    \node [anchor=west] at (0, -0.2) {\textbf{Tesis Conformes con Atribución Rigurosa (20\%) --- 3 trabajos}};
    \draw [fill=unmsm_gold, draw=unmsm_gray, line width=0.8pt] (0, -1.0) rectangle (2.4, -0.4);
    \node at (1.2, -0.7) [text=black, font=\small\sffamily\bfseries] {20\%};
\end{tikzpicture}
\caption*{\textbf{Figura 14.4.} \textit{Distribución de Omisiones de Fuentes en la Muestra de Trabajos de Grado de De Jesús G. (2016).} \scriptsize \textbf{Fuente:} Elaboración propia a partir de los datos cuantitativos del estudio citado.}
\end{figure}

No obstante la contundencia de estas cifras, el análisis epistemológico formulado por Palacios Carrasco et al. (2026) impone una indispensable \textbf{cautela metodológica}: una tasa del 80\% obtenida mediante una muestra intencional de 15 monografías no constituye bajo ningún concepto un parámetro epidemiológico generalizable a la totalidad del sistema universitario latinoamericano. La intención de De Jesús G. no fue inferir una prevalencia poblacional, sino desvelar una falla sistémica institucional: el hecho de que el 80\% de una muestra de tesis aprobadas exhibiera plagio severo demuestra la fragilidad de los procesos de arbitraje tutorial y la pasividad de los comités examinadores, quienes validaron trabajos sin realizar un cotejo riguroso de las fuentes invocadas.

\subsubsection{Prevalencia y Factores Actitudinales: El Estudio de Porto Castro (2024)}

Desde una perspectiva psicopedagógica y sociológica contemporánea, la investigación desarrollada por Porto Castro (2024) sobre una muestra de \textbf{310 estudiantes universitarios} en la Universidad de Santiago de Compostela (España) arroja luz sobre los determinantes comportamentales del fraude académico.

El estudio de Porto Castro demuestra empíricamente que la comisión de plagio rara vez responde a una vocación deliberada de engaño o dolo premeditado; en una alta proporción de los casos, obedece a una confluencia de factores contextuales:
\begin{itemize}
    \item \textbf{Gestión Ineficiente del Tiempo y Procrastinación Académica:} La postergación sistemática de las tareas de investigación genera picos extremos de ansiedad en los días previos a la fecha límite, induciendo al estudiante a recurrir al corte y confección apresurado de textos digitales.
    \item \textbf{Sobrecarga de Evaluaciones Heterogéneas:} La acumulación desarticulada de trabajos finales sin acompañamiento procesual incentiva la búsqueda de atajos cognitivos.
    \item \textbf{Analfabetismo Informacional y de Citación:} Desconocimiento operativo de las normativas de citación (normas APA, IEEE), ignorando que el parafraseo no citado constituye una infracción ética.
    \item \textbf{Percepción de Impunidad y Desafección Ética:} La convicción de que los docentes no leen con detenimiento los informes y que las herramientas de detección pueden burlarse fácilmente mediante alteraciones cosméticas.
\end{itemize}

\subsubsection{Correlaciones Estadísticas y Modelado del Fraude: Puche Villalobos (2025)}

En un trabajo reciente sobre el impacto de las tecnologías emergentes en el contexto universitario, Puche Villalobos (2025) explora las correlaciones entre el nivel de familiaridad digital, el uso acrítico de asistentes generativos y la frecuencia autoreportada de fraude académico. El autor informa una correlación de Pearson excepcionalmente elevada ($r = 0.980, p < 0.001$) entre el desconocimiento de las normativas de integridad institucional y la predisposición al uso deshonesto de tecnologías en trabajos de posgrado.

Frente a este dato, Palacios Carrasco et al. (2026) subrayan la necesidad de ejercer un riguroso discernimiento epistemológico: en ciencias exactas y aplicadas, es un axioma fundamental que \textit{correlación no demuestra causalidad} (\textit{cum hoc ergo propter hoc}). Una correlación cercana a la unidad ($r = 0.980$) en un estudio de percepción puede reflejar colinealidad instrumental, sesgos del cuestionario de autorreporte o el efecto distorsionador de variables confusoras intermedias (como la cultura de evaluación punitiva del centro universitario). Asumir linealmente que capacitar en normativa reducirá en un 98\% el fraude constituye una extrapolación ingenua: la transformación ética exige modificar los incentivos estructurales y los métodos de enseñanza de las facultades.

\begin{table}[htbp]
\centering
\small
\renewcommand{\arraystretch}{1.25}
\textbf{Tabla 14.12}\\[0.15cm]
\textit{Matriz Comparativa y Auditoría Metodológica de las Investigaciones Empíricas Base}\\[0.25cm]
\begin{tabularx}{\textwidth}{p{2.6cm}p{2.0cm}p{3.2cm}X p{3.0cm}}
\toprule
\textbf{Estudio y Fuente} & \textbf{Muestra ($N$)} & \textbf{Diseño y Contexto} & \textbf{Hallazgo Cuantitativo Central} & \textbf{Límites y Advertencia Metodológica} \\
\midrule
\textbf{De Jesús G. (2016)} & 15 tesis de grado (de 60 aprobadas) & Muestreo intencional documental; FACES-ULA (Venezuela). & 80\% de los trabajos (12/15) exhibieron omisiones graves de citas primarias y secundarias. & No extrapolable a toda la región; ilustra fallas en tutorías y comités de grado. \\
\textbf{Porto Castro (2024)} & 310 estudiantes universitarios & Diseño correlacional por encuesta; Santiago de Compostela (España). & Correlación significativa entre procrastinación, sobrecarga y recurrencia de plagio mosaico. & Datos basados en autorreporte estudiantil; sesgo potencial de deseabilidad social. \\
\textbf{Puche Villalobos (2025)} & Comunidad universitaria & Estudio exploratorio sobre fraude en la era de la IA generativa. & Correlación $r = 0.980$ entre desconocimiento normativo y conductas fraudulentas. & Correlación no prueba causalidad determinista; posibles variables confusoras. \\
\textbf{Corona \& Alonso (2026)} & Revisión sistemática PRISMA & Revisión de literatura en IES mexicanas sobre IA generativa. & Conceptualización del postplagio y obsolescencia del control punitivo tradicional. & Estudio cualitativo/normativo; requiere contrastación empírica longitudinal. \\
\bottomrule
\end{tabularx}
\vspace{0.15cm}
\raggedright \scriptsize \textbf{Nota.} Sistematización analítica formulada por el Grupo 9 (Palacios Carrasco et al., 2026).
\end{table}

\subsubsection{Las Cuatro Preguntas Críticas para Auditar Estudios sobre Integridad}

Para consolidar el pensamiento crítico en los estudiantes e investigadores de la UNMSM, se propone un esquema heurístico fundamentado en cuatro interrogantes metodológicas antes de asumir como válidas las conclusiones de cualquier artículo sobre plagio o fraude académico:

\begin{enumerate}
    \item \textbf{\textquestiondown Cuál es el tamaño de la muestra y cómo fueron seleccionados los participantes?} Verificar si se trata de un muestreo probabilístico riguroso o de un diseño por conveniencia/intencional, delimitando con precisión su validez externa.
    \item \textbf{\textquestiondown De qué contexto socioinstitucional y disciplinar provienen los datos?} Considerar las diferencias éticas, reglamentarias y laborales existentes entre diversas regiones y tradiciones universitarias.
    \item \textbf{\textquestiondown Qué instrumentos técnicos y metodológicos se emplearon para recolectar la información?} Distinguir los resultados obtenidos mediante cuestionarios de autorreporte subjetivo de aquellos fundados en auditorías ciegas de artefactos y repositorios reales.
    \item \textbf{\textquestiondown Se infiere una relación causal ilegítima a partir de una correlación estadística observada?} Recordar siempre que la covariación matemática entre dos magnitudes empíricas no demuestra un nexo de causa-efecto determinista.
\end{enumerate}

\subsection{Gobernanza Universitaria, Evaluación Formativa de Procesos y Debido Proceso}

\subsubsection{Los Límites Técnicos del Software de Similitud y los Detectores de IA}

En las últimas dos décadas, un número considerable de instituciones universitarias delegó erróneamente la tutela de la integridad académica en plataformas comerciales de análisis de coincidencias textuales (como Turnitin, SafeAssign o Urkund), bautizadas equívocamente como \textsl{\guillemotleft software antiplagio\guillemotright}. Esta delegación técnica generó una distorsión pedagógica severa: la creencia de que un algoritmo computacional es capaz de dictaminar la presencia o ausencia de plagio moral.

Como demuestran Palacios Carrasco et al. (2026), estas plataformas operan mediante algoritmos deterministas de búsqueda de subcadenas comunes ($k$-gramas y distancia de Levenshtein) contra bases de datos indexadas. Su único producto técnico es un \textbf{reporte cuantitativo de coincidencias porcentuales}, el cual resulta ciego a la semántica y a la naturaleza del discurso científico:
\begin{itemize}
    \item \textbf{Falsos Positivos Inevitables:} El software resalta mecánicamente como supuestas coincidencias las listas bibliográficas estructuradas según normativas obligatorias (APA, IEEE), las fórmulas matemáticas estándar (como la integral gaussiana de Poisson), la carátula institucional reglamentaria de la universidad y la fraseología metodológica recurrente (\textsl{\guillemotleft para contrastar la hipótesis nula se aplicó una prueba t de Student\guillemotright}).
    \item \textbf{Falsos Negativos Críticos:} Las herramientas tradicionales son incapaces de advertir el plagio mosaico con sustitución hábil de sinónimos, la traducción translingüística de artículos en inglés o ruso, la apropiación de ideas no publicadas y el texto generado sintéticamente por modelos fundacionales de IA.
    \item \textbf{Inestabilidad Radical de los Detectores de Texto Sintético:} Los detectores de IA generativa (basados en métricas de perplejidad y variabilidad léxica o \textit{burstiness}) presentan tasas de falsos positivos inaceptables, las cuales se disparan dramáticamente contra estudiantes cuya lengua materna no es el inglés o que redactan con estilos formales muy estructurados.
\end{itemize}

La ineficacia de confiar la integridad al software ha sido demostrada empíricamente en rigurosos estudios internacionales. Foltýnek et al. (2020) evaluaron decenas de sistemas comerciales de detección de similitud, constatando que no logran identificar paráfrasis elaboradas ni traducciones. Asimismo, Weber-Wulff et al. (2023) y la red ENAI (Foltýnek et al., 2023) evaluaron detectores de texto generado por IA, concluyendo que ninguno ofrece la precisión necesaria para fundamentar acusaciones académicas. Por consiguiente, especialistas como Phillip Dawson (2021) y James M. Lang (2013) postulan que la seguridad evaluativa universitaria exige abandonar la vigilancia técnica punitiva y priorizar el diseño de tareas auténticas y evaluaciones formativas procesuales.

En consecuencia, la universidad debe consagrar como principio innegociable que \textbf{un porcentaje de coincidencia no prueba ni descarta absolutamente nada por sí solo}. El reporte computacional constituye un simple indicio preliminar que exige indefectiblemente la \textbf{revisión humana experta y contextualizada} del docente o del comité disciplinar, analizando la intencionalidad, el aporte conceptual y la veracidad de las fuentes citadas.

\subsubsection{Rediseño Curricular: De la Evaluación de Productos a la Evaluación de Procesos}

El fracaso demostrado del paradigma punitivo-reactivo fundamentado en la vigilancia por software impone una transición pedagógica ineludible: la adopción de la \textbf{evaluación formativa de procesos}. Cuando la evaluación académica se reduce a exigir un único informe final en la semana 16 del ciclo académico, se crean las condiciones perfectas para el fraude: el docente solo ve el producto culminado y el estudiante procrastina hasta encontrarse frente a una presión temporal inmanejable.

La estrategia propuesta por el Grupo 9 (Palacios Carrasco et al., 2026) consiste en fragmentar secuencialmente la investigación a lo largo de las 16 semanas del semestre académico en cuatro hitos intermedios articulados:
\begin{enumerate}
    \item \textbf{Hito 1 (Semana 4):} Planteamiento de la pregunta de investigación, justificación y matriz preliminar de variables.
    \item \textbf{Hito 2 (Semana 8):} Entrega de la Bitácora de Lectura Científica con un mínimo de diez fuentes primarias arbitradas verificadas mediante DOI.
    \item \textbf{Hito 3 (Semana 12):} Entrega del borrador integral de trabajo. El docente proporciona retroalimentación formativa cualitativa sobre la coherencia del modelo y la calidad del parafraseo, permitiendo la corrección sin penalización.
    \item \textbf{Hito 4 (Semana 16):} Entrega del informe final acompañado de la matriz CRediT, el apéndice de Declaración de IA y la \textbf{sustentación oral individual obligatoria} en pizarra o revisión de código fuente.
\end{enumerate}

\begin{table}[htbp]
\centering
\small
\renewcommand{\arraystretch}{1.2}
\textbf{Tabla 14.13}\\[0.15cm]
\textit{Cronograma Semestral de 16 Semanas para la Gestión Procesual de la Integridad Académica}\\[0.25cm]
\resizebox{\textwidth}{!}{%
\begin{tabularx}{\textwidth}{p{1.8cm}p{4.2cm}X p{3.2cm}}
\toprule
\textbf{Semana} & \textbf{Hito Académico Exigido} & \textbf{Evidencia Trazable en el Repositorio} & \textbf{Mecanismo Docente de Supervisión} \\
\midrule
\textbf{Semana 1} & Presentación del sílabo y normas de integridad académica e IA. & Firma digital del Compromiso de Integridad y Ética Universitaria. & Registro en el aula virtual; resolución de dudas normativas. \\
\textbf{Semana 4} & Hito 1: Pregunta de investigación, hipótesis y diseño del modelo. & Documento de especificación de requisitos y delimitación matemática. & Revisión y retroalimentación de la viabilidad conceptual. \\
\textbf{Semana 8} & Hito 2: Bitácora de Lectura Científica y curación de datos. & Tabla estructurada con 10 fuentes primarias, DOIs verificados y paráfrasis. & Auditoría cruzada de DOIs y control de fuentes en Crossref. \\
\textbf{Semana 12} & Hito 3: Primer borrador integral y código preliminar del modelo. & Repositorio git con historial de commits y borrador en formato \LaTeX. & Revisión formativa cualitativa; corrección temprana de errores. \\
\textbf{Semana 15} & Hito 4: Informe final refinado, matriz CRediT y declaración de IA. & Manuscrito final, repositorio reproducible y plantilla de transparencia. & Análisis contextual del reporte de similitud y verificación. \\
\textbf{Semana 16} & Defensa oral pública y examen individual de código en pizarra. & Registro de preguntas orales y evaluación de solvencia cognitiva. & Calificación definitiva basada en la rúbrica de proceso. \\
\bottomrule
\end{tabularx}%
}
\vspace{0.15cm}
\raggedright \scriptsize \textbf{Nota.} Cronograma formativo formulado por Palacios Carrasco et al. (2026).
\end{table}

\subsubsection{Rúbrica Analítica Multidimensional de Evaluación de Integridad}

Para dotar de objetividad y equidad a la evaluación universitaria, se estructura una rúbrica analítica que pondera cinco dimensiones fundamentales a través de tres niveles progresivos de dominio: Inicial, En Proceso y Logrado.

\begin{table}[htbp]
\centering
\small
\renewcommand{\arraystretch}{1.25}
\textbf{Tabla 14.14}\\[0.15cm]
\textit{Rúbrica Analítica Multidimensional para la Evaluación Formativa de Procesos de Integridad}\\[0.25cm]
\begin{tabularx}{\textwidth}{p{2.8cm}X X X}
\toprule
\textbf{Criterio Evaluativo} & \textbf{Nivel Inicial (Insuficiente)} & \textbf{Nivel En Proceso (Aceptable)} & \textbf{Nivel Logrado (Sobresaliente)} \\
\midrule
\textbf{Trazabilidad de Fuentes y Citas} & Fuentes no académicas; citas omitidas o erróneas; ausencia de DOIs en la lista final. & Fuentes académicas pero con inconsistencias menores en formato APA o enlaces caídos. & Todas las fuentes son primarias y arbitradas; DOIs verificados y citación formal impecable. \\
\textbf{Calidad de Paráfrasis y Síntesis} & Predominio de plagio mosaico, copia servil o sustitución superficial de sinónimos. & Paráfrasis correcta en la mayoría de pasajes pero con dependencia de la sintaxis original. & Reconstrucción conceptual original, fluida y con voz autoral propia; ideas ajenas bien integradas. \\
\textbf{Reproducibilidad y Código} & Código huérfano sin entorno ni semillas; datasets sin procesar no disponibles ni citados. & Código funcional pero sin fijación determinista de semillas o con librerías sin versionado. & Repositorio git reproducible al 100\%; \texttt{requirements.txt}, semillas fijadas y datos citados. \\
\textbf{Matriz CRediT y Coautoría} & Asignación simétrica ficticia sin correspondencia con los aportes reales del equipo. & Matriz presente pero con descripciones vagas o genéricas de los roles de los coautores. & Roles CRediT perfectamente delimitados y concordantes con la auditoría de \textit{commits}. \\
\textbf{Transparencia en Uso de IA} & Uso encubierto no declarado; o entrega de declaraciones genéricas incompletas. & Declara las herramientas empleadas pero omite los \textit{prompts} o la verificación empírica. & Declaración completa (herramienta, prompts, cambios y verificación empírica de datos). \\
\textbf{Solvencia en Defensa Oral} & Incapacidad para explicar las ecuaciones, el código o responder por lo firmado. & Explica la lógica general pero titubea al justificar decisiones matemáticas complejas. & Defiende con absoluta solvencia teórica y técnica cada decisión y línea de código suscrita. \\
\bottomrule
\end{tabularx}
\vspace{0.15cm}
\raggedright \scriptsize \textbf{Nota.} Rúbrica formativa diseñada por Grupo 9 (Palacios Carrasco et al., 2026).
\end{table}

\subsubsection{Cláusulas Institucionales Modelo para Sílabos Universitarios}

La prevención eficaz exige que las reglas de juego éticas queden perfectamente explícitas desde el primer día de clases. A continuación se presentan las cláusulas modelo diseñadas para su incorporación en los sílabos de la Facultad de Ciencias Matemáticas:

\begin{table}[htbp]
\centering
\small
\renewcommand{\arraystretch}{1.25}
\textbf{Tabla 14.15}\\[0.15cm]
\textit{Cláusulas Modelo de Integridad Académica y Gobernanza de IA para Sílabos}\\[0.25cm]
\begin{tabularx}{\textwidth}{p{3.0cm}X}
\toprule
\textbf{Cláusula Sílabica} & \textbf{Texto Normativo Sugerido para Adopción Institucional en UNMSM} \\
\midrule
\textbf{1. Compromiso de Integridad y Firma Solidaria} & \textsl{\guillemotleft La suscripción de cualquier trabajo, examen o proyecto en este curso confiere a los firmantes la condición de coautores responsables solidarios. Toda afirmación, dato numérico, algoritmo o fragmento discursivo tomado de fuentes ajenas debe ser citado con estricta sujeción a las normas APA 7.ª edición. La firma implica la capacidad irrenunciable de defender oralmente la totalidad de la obra\guillemotright}. \\
\textbf{2. Política de Uso Transparente de IA Generativa} & \textsl{\guillemotleft Se permite el uso de modelos fundacionales de IA exclusivamente con propósitos formativos de andamiaje cognitivo (lluvia de ideas, explicación de dudas conceptuales y depuración sintáctica de código propio). Queda estrictamente prohibida la subcontratación algorítmica consistente en generar párrafos o trabajos completos sin autoría humana. Todo uso debe declararse en el apéndice metodológico oficial\guillemotright}. \\
\textbf{3. Verificación Empírica Obligatoria de Fuentes} & \textsl{\guillemotleft Queda terminantemente prohibido incorporar citas bibliográficas generadas por IA que no hayan sido verificadas empíricamente por el estudiante en bases de datos indexadas. La inclusión de referencias apócrifas o DOIs inexistentes se considerará falta grave a la veracidad científica, independientemente de la intencionalidad del autor\guillemotright}. \\
\textbf{4. Garantía de Debido Proceso y Defensa Oral} & \textsl{\guillemotleft Ningún estudiante será sancionado basándose exclusivamente en porcentajes emitidos por software de similitud o detectores automáticos. Ante cualquier discordancia o sospecha de fraude, se activará el protocolo de debido proceso, requiriendo las evidencias de la bitácora y celebrando una audiencia oral para escuchar los descargos del estudiante\guillemotright}. \\
\bottomrule
\end{tabularx}
\vspace{0.15cm}
\raggedright \scriptsize \textbf{Nota.} Formulación institucional propuesta por Palacios Carrasco et al. (2026).
\end{table}

\subsubsection{Protocolo Institucional de Debido Proceso ante Sospechas de Fraude}

El Estado constitucional de derecho y la autonomía universitaria imponen que ninguna sanción académica sea impuesta de forma arbitraria, sumaria o desproporcionada. Cuando un docente o comité detecta anomalías severas en un trabajo, debe ejecutarse el siguiente protocolo estructurado en seis fases sucesivas:

\begin{figure}[htbp]
\centering
\begin{tikzpicture}[
    node distance=1.0cm and 0.5cm,
    font=\footnotesize\sffamily,
    stage/.style={rectangle, draw=unmsm_red, line width=0.8pt, fill=bg_marfil, rounded corners=2pt, text width=2.1cm, text centered, minimum height=1.4cm},
    arrow/.style={-{Latex[length=2.5mm]}, thick, draw=unmsm_darkred}
]
    \node [stage] (s1) {\textbf{Fase 1}\\[0.1cm]Auditoría Contextual};
    \node [stage, right=of s1] (s2) {\textbf{Fase 2}\\[0.1cm]Tipificación de la Falta};
    \node [stage, right=of s2] (s3) {\textbf{Fase 3}\\[0.1cm]Solicitud de Evidencias};
    \node [stage, right=of s3] (s4) {\textbf{Fase 4}\\[0.1cm]Audiencia de Descargos};
    \node [stage, below=1.0cm of s4] (s5) {\textbf{Fase 5}\\[0.1cm]Dictamen Motivado};
    \node [stage, left=of s5] (s6) {\textbf{Fase 6}\\[0.1cm]Medida Proporcional};
    
    \draw [arrow] (s1) -- (s2);
    \draw [arrow] (s2) -- (s3);
    \draw [arrow] (s3) -- (s4);
    \draw [arrow] (s4) -- (s5);
    \draw [arrow] (s5) -- (s6);
\end{tikzpicture}
\caption*{\textbf{Figura 14.5.} \textit{Diagrama de Flujo del Protocolo de Debido Proceso ante Sospechas de Fraude Académico.} \scriptsize \textbf{Fuente:} Adaptado de Palacios Carrasco et al. (2026).}
\end{figure}

\begin{enumerate}
    \item \textbf{Fase 1: Auditoría Contextual del Reporte:} El docente examina en detalle las coincidencias resaltadas por el software, descartando manualmente referencias, carátulas y fórmulas comunes.
    \item \textbf{Fase 2: Tipificación Preliminar de la Presunta Falta:} Se clasifica formalmente la naturaleza de la irregularidad observada: \textsl{¿se trata de un error metodológico de citación involuntario, un plagio mosaico por mala técnica de paráfrasis o una subcontratación fraudulenta deliberada?}
    \item \textbf{Fase 3: Requerimiento de Evidencias de Bitácora:} Se solicita al estudiante exhibir su Bitácora de Lectura Científica, el historial de \textit{commits} del repositorio git y las notas manuscritas que demuestren la evolución cronológica del trabajo.
    \item \textbf{Fase 4: Audiencia de Descargo y Diálogo Académico:} Se convoca al estudiante a una entrevista oral individual en la que se le solicita explicar la lógica matemática del pasaje cuestionado y justificar sus decisiones metodológicas.
    \item \textbf{Fase 5: Dictamen Motivado del Comité o Docente:} Se emite un acta formal por escrito donde se exponen con claridad los fundamentos racionales y las evidencias cotejadas que sustentan la resolución adoptada.
    \item \textbf{Fase 6: Aplicación de Medidas Pedagógicas y Disciplinarias Proporcionadas:} Si se concluye que hubo desconocimiento metodológico sin dolo, la medida debe ser formativa (rehacer el apartado bajo tutoría y asistir a un taller de citación); si se comprueba dolo manifiesto o falsificación consciente de datos, se aplican las sanciones disciplinarias graduadas según el estatuto universitario (desaprobación de la asignatura o apertura de proceso administrativo).
\end{enumerate}

\subsubsection{Matriz de Responsabilidad Compartida e Indicadores de Madurez Ética}

La construcción de una cultura de integridad no puede recaer de manera unilateral sobre los hombros del estudiante ni transformarse en una carga policial para el docente; constituye una empresa colectiva que compromete a los cuatro estamentos de la comunidad universitaria.

\begin{table}[htbp]
\centering
\small
\renewcommand{\arraystretch}{1.25}
\textbf{Tabla 14.16}\\[0.15cm]
\textit{Matriz de Responsabilidad Compartida en la Universidad Nacional Mayor de San Marcos}\\[0.25cm]
\begin{tabularx}{\textwidth}{p{3.2cm}X}
\toprule
\textbf{Estamento Universitario} & \textbf{Compromisos Deontológicos y Acciones Operativas Ineludibles} \\
\midrule
\textbf{Estudiantes} & Llevar bitácoras de lectura rigurosas; citar con exactitud matemática; verificar empíricamente cada fuente; declarar con honestidad el uso de IA; y estar en condiciones de defender oralmente cada línea suscrita. \\
\textbf{Docentes} & Explicitar las reglas de integridad desde el sílabo; fragmentar las entregas semestrales; retroalimentar borradores intermedios; evitar la dependencia ciega del software; y evaluar mediante explicaciones orales. \\
\textbf{Departamentos y Escuelas} & Capacitar de forma continua al profesorado en el uso crítico de herramientas de análisis; estructurar comités de ética ágiles; y estandarizar las rúbricas analíticas de evaluación de procesos. \\
\textbf{Autoridades Universitarias} & Promover políticas institucionales formativas; adecuar los reglamentos a la era de la IA; dotar de repositorios institucionales abiertos; y garantizar el debido proceso y la defensa de la comunidad. \\
\bottomrule
\end{tabularx}
\vspace{0.15cm}
\raggedright \scriptsize \textbf{Nota.} Elaboración propia a partir del marco institucional de Palacios Carrasco et al. (2026).
\end{table}

Para evaluar si la política institucional de integridad está madurando favorablemente, los departamentos académicos no deben limitarse a contar cuántos estudiantes son castigados al año. Por el contrario, deben monitorear indicadores cuantitativos y cualitativos de salud académica:
\begin{itemize}
    \item \textbf{Tasa de Cumplimiento de Hitos Intermedios:} Proporción de estudiantes que entregan oportunamente la propuesta inicial, la bitácora anotada y el borrador de trabajo.
    \item \textbf{Índice de Calidad de Paráfrasis:} Evaluación muestral de la divergencia sintáctica y originalidad argumentativa en los borradores de tesis de grado.
    \item \textbf{Trazabilidad de Declaraciones de IA:} Porcentaje de informes que incluyen transparentemente la plantilla de declaración y verificación de prompts.
    \item \textbf{Resolución Formativa de Incidentes:} Número de casos de similitud resueltos constructivamente mediante tutoría y reescritura, frente a procesos que derivan en sanciones sumarias punitivas.
\end{itemize}

\subsubsection{Escenarios Críticos y Dilemas Éticos para el Debate en Aula}

Con la finalidad de ejercitar el juicio ético y la casuística deontológica en las aulas de la UNMSM, se presentan cinco escenarios deliberativos inspirados en la praxis cotidiana de la investigación computacional:

\begin{enumerate}
    \item \textbf{Escenario 1 (El Asesor Pasivo):} Un equipo de tesistas culmina una investigación de alto impacto y, por agradecimiento al docente que les facilitó el acceso al laboratorio de cómputo (pero que no participó en el diseño ni leyó el manuscrito), deciden incluirlo como primer coautor. \textit{\textquestiondown Qué tipología de autoría ficticia se configura? \textquestiondown Cómo debió reconocerse éticamente su valioso apoyo logístico?}
    \item \textbf{Escenario 2 (La Paráfrasis no Citada):} Un estudiante lee con atención una definición matemática propuesta en una tesis europea, la redacta íntegramente con sus propias palabras técnicas y no incluye ninguna cita parentética, alegando que \textsl{\guillemotleft las palabras son totalmente mías y por ende el párrafo me pertenece\guillemotright}. \textit{\textquestiondown Constituye plagio esta conducta? \textquestiondown Qué dimensión ética y jurídica fue vulnerada?}
    \item \textbf{Escenario 3 (La Reutilización de Código sin Licencia):} Un estudiante de computación científica extrae una función de optimización de un repositorio público en GitHub con licencia GNU GPLv3, la inserta en un proyecto comercial universitario cerrado sin declarar su origen ni incluir la licencia matriz. \textit{\textquestiondown Qué principios éticos y legales se infringieron? \textquestiondown Cómo se subsana formalmente la omisión?}
    \item \textbf{Escenario 4 (La Traducción Algorítmica Encubierta):} Una estudiante solicita a un modelo de IA generativa que traduzca un artículo científico publicado en ruso, ajusta la redacción y lo presenta como un marco teórico propio sin mencionar la obra original ni la herramienta traductora. \textit{\textquestiondown Cómo se categoriza esta práctica ante un comité disciplinar?}
    \item \textbf{Escenario 5 (El Error Metodológico Descubierto):} Tres días después de haber entregado el informe final calificado con nota máxima, un estudiante descubre un fallo grave en la semilla aleatoria que invalida la convergencia del algoritmo principal informado. \textit{\textquestiondown Cuál es la conducta deontológica que el valor de la valentía académica impone ejecutar de inmediato?}
\end{enumerate}

\subsubsection{Glosario Terminológico Institucional}

\begin{itemize}
    \item \textbf{Alucinación Algorítmica:} Salida espuria generada por un modelo probabilístico de lenguaje en la que se formulan datos, hechos, demostraciones o referencias bibliográficas inexistentes con una apariencia de absoluta verosimilitud formal.
    \item \textbf{Autoplagio (Reciclaje Textual):} Práctica éticamente reprobable consistente en presentar como nuevo e inédito un trabajo, datos o texto propio que ya fue previamente calificado, publicado o evaluado, sin declarar dicha reutilización.
    \item \textbf{Autoría Fantasma (\textit{Ghost Authorship}):} Omisión deliberada del reconocimiento autoral de quien contribuyó sustantivamente a la investigación; comprende también la utilización encubierta de IA generativa para redactar sustancialmente la obra.
    \item \textbf{Bitácora de Lectura Científica:} Instrumento metodológico de registro inmutable donde el investigador documenta los metadatos, citas textuales, paráfrasis y valoración crítica de cada fuente consultada antes de redactar.
    \item \textbf{Capa de Aseguramiento (\textit{Assurance Layer}):} Conjunto estructurado de mecanismos de verificación humana (defensas orales, auditoría SAST, declaración de IA) que garantiza la calidad y veracidad del trabajo académico.
    \item \textbf{Convenio de Berna:} Tratado internacional fundacional (1886) que establece el principio de trato nacional y la protección automática transfronteriza de los derechos morales y patrimoniales de los creadores de obras literarias y artísticas.
    \item \textbf{Derecho de Cita:} Excepción legal formal (Art. 44 inc. a del D.L. 822) que autoriza la reproducción no lucrativa de breves extractos de obras protegidas con fines docentes o de investigación, exigiendo atribución y usos honrados.
    \item \textbf{Derechos Morales:} Facultades perpetuas, inalienables e inembargables que vinculan al autor humano con su obra intelectual, destacando los derechos de paternidad e integridad.
    \item \textbf{Derechos Patrimoniales:} Prerrogativas de aprovechamiento económico exclusivo sobre la obra (reproducción, comunicación pública, distribución, traducción), susceptibles de cesión temporal a terceros.
    \item \textbf{Fraude por Subcontratación Algorítmica:} Práctica fraudulenta consistente en delegar en un sistema de IA la redacción de tareas o la resolución de problemas para presentarlos como producto de la cognición del estudiante.
    \item \textbf{Integridad Académica:} Compromiso activo y transversal con los valores de honestidad, confianza, justicia, respeto, responsabilidad y valentía en la producción, docencia y difusión del saber científico.
    \item \textbf{Plagio Mosaico (\textit{Patchwork Plagiarism}):} Reformulación cosmética consistente en sustituir sinónimos o alterar el orden de cláusulas conservando el entramado lógico y argumentativo de la fuente ajena sin atribuirla.
    \item \textbf{Postplagio:} Teoría contemporánea que reconoce la hibridación de la escritura humana y artificial, desplazando el control punitivo de coincidencias textuales hacia la evaluación de la agencia cognitiva y la verificación ética.
    \item \textbf{Reproducibilidad Científica:} Capacidad técnica de un par académico independiente de recrear con precisión exacta los resultados numéricos informados en una publicación a partir de los datos brutos y el código fuente.
    \item \textbf{Taxonomía CRediT:} Estándar internacional normalizado (ANSI/NISO Z39.104-2022) que desglosa en 14 roles específicos las contribuciones intelectuales, metodológicas y técnicas de los colaboradores de un proyecto científico.
    \item \textbf{Usos Honrados (\textit{Fair Use}):} Doctrina jurídica que pondera la licitud de la utilización no autorizada de una obra protegida atendiendo al propósito del uso, naturaleza de la obra, cantidad extraída y efecto sobre el mercado.
\end{itemize}

% ==============================================================================
% NOTAS BIBLIOGRÁFICAS E HISTÓRICAS DEL CAPÍTULO 11
% ==============================================================================
\newpage
\markboth{CAPÍTULO 14\quad NOTAS BIBLIOGRÁFICAS E HISTÓRICAS}{}
\vspace*{0.8cm}
\noindent{\Huge \textbf{Notas bibliográficas\\[0.2cm] e históricas}}\\[0.8cm]

\textbf{Antecedentes Teóricos y Normativos Fundacionales:}\\[0.2cm]
La protección formal del intelecto y la demarcación del plagio poseen un itinerario histórico milenario. Desde la célebre censura del poeta romano Marcial en el siglo I d.C. ---quien acuñó metafóricamente el vocablo latino \textit{plagium} (el secuestro de esclavos ajenos) para increpar a quienes recitaban sus versos como propios---, hasta la promulgación del \textit{Estatuto de la Reina Ana} en Inglaterra (1710), la civilización occidental transitó progresivamente desde un régimen de privilegios reales monárquicos hacia el reconocimiento formal del derecho del creador sobre su obra.

El hito normativo de mayor trascendencia global fue la adopción del \textbf{Convenio de Berna para la Protección de las Obras Literarias y Artísticas (1886)}, el cual consagró por primera vez en el derecho internacional el principio de protección automática (sin necesidad de registro formal previo) y blindó los derechos morales de paternidad e integridad como salvaguarda innegociable de la personalidad del autor frente a las presiones del mercado editorial transfronterizo.

En la esfera académica, la teorización moderna de la integridad como disciplina axiológica independiente floreció a partir de los trabajos pioneros de Donald McCabe, quien en 1992 fundó el \textbf{Center for Academic Integrity} (evolucionado hoy en el ICAI). McCabe demostró empíricamente que los códigos de honor tradicionales y la vigilancia policíaca resultaban ineficaces si no estaban respaldados por una cultura comunitaria de confianza mutua, justicia y valentía moral. Asimismo, la entrada en vigencia del \textbf{Decreto Legislativo N.º 822 (1996)} en el Perú modernizó el marco legal nacional, adaptándolo a las exigencias de la Organización Mundial del Comercio (OMC) y fijando un estricto equilibrio entre los derechos exclusivos del autor y las excepciones de interés público para la enseñanza universitaria.\\[0.4cm]

\textbf{Contribuciones de las Investigaciones Base del Capítulo:}\\[0.2cm]
El aparato conceptual, doctrinal y empírico desarrollado a lo largo de este capítulo se fundamenta en la integración rigurosa de las investigaciones científicas aportadas por el Grupo 9 de la UNMSM y la literatura internacional contemporánea:
\begin{itemize}
    \item \textbf{International Center for Academic Integrity [ICAI] (2021):} Provee la matriz axiológica fundacional articulada en los seis valores universales (honestidad, confianza, justicia, respeto, responsabilidad y valentía), sirviendo de base para la reformulación formativa de los reglamentos universitarios frente al punitivismo estéril.
    \item \textbf{De Jesús G., M. I. (2016) [Universidad de Los Andes, Venezuela]:} Aporta una investigación referencial indispensable en el contexto hispanoamericano, auditando empíricamente la prevalencia del 80\% de omisiones documentales en tesis aprobadas y tipificando con rigor las patologías de las autorías por regalo, intercambio y estatus.
    \item \textbf{Porto Castro, A. M. (2024) [Universidad de Santiago de Compostela / Editorial Redipe]:} Fundamenta la sociología y psicología del fraude estudiantil mediante su investigación empírica sobre 310 universitarios, demostrando la correlación directa entre sobrecarga académica, procrastinación y plagio mosaico.
    \item \textbf{Corona Domínguez, B., y Alonso Nuñez, M. C. (2026) [Ciencia y Reflexión]:} Desarrollan la revisión sistemática que introduce formalmente en la región el paradigma del \textit{postplagio}, sustentando la obsolescencia de los detectores automáticos y la necesidad de priorizar la agencia cognitiva y la autoría híbrida responsable.
    \item \textbf{Puche Villalobos, D. J. (2025) [Revista Digital de Investigación y Postgrado]:} Provee el modelado cuantitativo de correlaciones ($r = 0.980$) sobre la interacción entre desconocimiento normativo y fraude tecnológico, abriendo el debate epistemológico sobre la diferenciación indispensable entre covariación estadística y causalidad.
    \item \textbf{Miao, F., y Holmes, W. (2023) [UNESCO]:} Formulan la directriz política global más influyente de la actualidad sobre el uso ético de la IA en la educación y la ciencia, consagrando el principio rector vinculante de la \textbf{primacía de la agencia humana} y la prohibición de atribuir coautorías a sistemas computacionales.
    \item \textbf{American Psychological Association [APA] (2022):} Sistematiza las directrices internacionales sobre autoplagio, citas parentéticas rigurosas y autorizaciones formales requeridas para la adaptación de figuras y tablas complejas bajo normativas de propiedad intelectual.
\end{itemize}
\vspace{0.4cm}

\textbf{Lecturas Sugeridas Comentadas:}\\[0.2cm]
\begin{enumerate}
    \item \textbf{Eaton, S. E. (2021):} \textit{Plagiarism in Higher Education: Tackling Tough Topics in Academic Integrity}. Una obra magistral que anticipó el advenimiento del postplagio y propone estrategias curriculares concretas para transformar la cultura evaluativa institucional.
    \item \textbf{Brand, A., Allen, L., Altman, M., Hlava, M., \& Scott, J. (2015):} \textit{Beyond authorship: attribution, contribution, collaboration, and credit}. El documento fundacional que estructuró la Taxonomía CRediT, lectura imprescindible para gestionar la coautoría en laboratorios de computación y ciencias exactas.
    \item \textbf{UNESCO (2021):} \textit{Recomendación sobre la Ética de la Inteligencia Artificial}. El marco normativo multilateral universal que orienta el desarrollo de tecnologías inteligentes al servicio de la dignidad humana, los derechos fundamentales y la justicia social.
\end{enumerate}

% ==============================================================================
% RESUMEN DEL CAPÍTULO 11
% ==============================================================================
\newpage
\markboth{CAPÍTULO 14\quad RESUMEN}{}
\vspace*{0.6cm}
{\Huge \textbf{Resumen del capítulo 14}}
\vspace{0.5cm}

La salvaguarda de la integridad académica, el respeto irrestricto al derecho de autor y la navegación responsable del desafío epistémico que plantea la Inteligencia Artificial Generativa exigen abandonar el punitivismo reactivo y articular un diseño pedagógico transparente que haga visible el proceso cognitivo completo del estudiante y asiente una cultura de responsabilidad compartida.

\vspace{0.3cm}

\begin{resumenlist}
    \item \textbf{Fundamentos Axiológicos y Dualidad Jurídica:} La ética académica y el derecho de autor (D.L. 822) operan como regímenes complementarios pero distintos: la ley positiva tutela la forma exterior distinguiendo derechos morales de patrimoniales, mientras la ética universitaria protege ideas y métodos. Los seis valores del ICAI operativizan las decisiones cotidianas de investigación.
    \item \textbf{Taxonomías Morfológicas y Justicia Autorial:} El plagio abarca la copia servil, el plagio mosaico (parafraseo cosmético) y el autoplagio no declarado. Su prevención exige dominar la paráfrasis legítima y la bitácora de lectura. Asimismo, erradicar autorías ficticias (por regalo o intercambio) demanda aplicar los criterios ICMJE y la Taxonomía CRediT.
    \item \textbf{Inteligencia Artificial, Postplagio y Verificación:} La disrupción de los LLMs quiebra la identidad entre prosa coherente y asimilación conceptual. La universidad debe delimitar el uso formativo (andamiaje cognitivo) del sustitutivo (subcontratación fraudulenta). Frente a las alucinaciones algorítmicas, el postplagio impone la primacía de la agencia humana y la verificación empírica.
    \item \textbf{Gobernanza Institucional y Evaluación Formativa:} Los reportes de coincidencia son indicios que exigen revisión humana experta. La transformación estructural del aula universitaria requiere sustituir la entrega final por un cronograma de 16 semanas con hitos procesuales trazables, garantizando un protocolo formal de debido proceso y medidas pedagógicas proporcionales.
\end{resumenlist}

\vspace{0.4cm}
\begin{cajabeige}
\textbf{Síntesis:} La integridad universitaria no resulta de la vigilancia pasiva por software, sino de un diseño pedagógico transparente: reglas claras en el sílabo, procesos con trazabilidad auditable, verificación empírica de fuentes y un compromiso ético solidario entre estudiantes, docentes y autoridades institucionales.
\end{cajabeige}

% ==============================================================================
% CONCLUSIONES DEL CAPÍTULO 11
% ==============================================================================
\newpage
\markboth{CAPÍTULO 14\quad CONCLUSIONES}{}
\vspace*{0.6cm}
{\Huge \textbf{Conclusiones del capítulo 14}}
\vspace{0.5cm}

\subsubsection*{Conclusión principal}
La universidad contemporánea sólo podrá integrar de manera fecunda las herramientas digitales y los modelos de Inteligencia Artificial Generativa si abandona la ilusión estéril de una respuesta exclusivamente punitiva o tecnocrática y asume que la integridad académica reside en el desarrollo demostrable de la agencia cognitiva y la rendición de cuentas del investigador. Ninguna tecnología puede sustituir el acto reflexivo humano de comprender los fundamentos matemáticos de un modelo, cotejar la veracidad histórica de una fuente o hacerse responsable moral y jurídicamente por las consecuencias del conocimiento suscrito. La autenticidad en la era de los algoritmos no se define por la ausencia de máquinas auxiliares, sino por la soberanía epistemológica con la que el ser humano dirige, audita y valida la creación del saber científico.

\subsubsection*{Aportes y límites}
\textbf{Aportes:} Este capítulo monográfico sistematiza de forma orgánica la doble vertiente jurídica (D.L. 822 de Perú y Convenio de Berna) y axiológica (ICAI 2021) de la propiedad intelectual; desmitifica el funcionamiento algorítmico del software de similitud y los detectores de texto sintético; formula herramientas operativas listas para adopción inmediata (la Bitácora de Lectura Científica, el Ciclo de Trabajo Íntegro de 7 fases, la Declaración de Uso de IA en 5 dimensiones y el Árbol de Decisión); y propone un modelo curricular integral de 16 semanas fundado en la fragmentación de entregas, la Taxonomía CRediT, rúbricas analíticas de proceso y protocolos institucionales de debido proceso.\\[0.15cm]
\textbf{Límites:} La rápida evolución técnica de las arquitecturas de modelos fundacionales (sistemas multiagente, modelos multimodales y agentes autónomos de código) supera con creces la velocidad de adaptación de las reglamentaciones jurídicas y estatutos universitarios, lo que demanda revisiones continuas de las directrices sílabicas. Asimismo, los hallazgos empíricos cuantitativos revisados en la literatura (como las tasas de omisión de De Jesús G. o las correlaciones de Puche Villalobos) provienen de contextos y diseños específicos que no pueden extrapolarse como leyes universales aplicables idénticamente a todas las disciplinas de las ciencias exactas.

\subsubsection*{Preguntas abiertas}
\textquestiondown De qué manera deberán reconfigurarse estructuralmente los instrumentos normativos y los exámenes de grado y posgrado en la educación superior para certificar la originalidad y la solvencia cognitiva de un futuro científico en un horizonte donde los agentes autónomos de inteligencia artificial no operen como herramientas externas de consulta, sino como sistemas co-creadores omnipresentes integrados de forma nativa e invisible en los entornos de computación y redacción de la ciencia global?

% ==============================================================================
% REFERENCIAS BIBLIOGRÁFICAS DEL CAPÍTULO 11 (NORMAS APA 7)
% ==============================================================================
\newpage
\markboth{CAPÍTULO 14\quad REFERENCIAS}{}
\vspace*{0.8cm}
{\Huge \textbf{Referencias del Capítulo 14}}\\[0.8cm]

\begin{itemize}[leftmargin=1.5em, label={}, itemsep=0.22ex, parsep=0pt]
\small\raggedright
\interlinepenalty=10000
\item American Psychological Association. (2020). \textit{Publication manual of the American Psychological Association} (7.ª ed.). American Psychological Association. \url{https://doi.org/10.1037/0000165-000}
\item American Psychological Association. (2022). \textit{Plagiarism and academic integrity guidelines: Direct quotes, paraphrasing, and copyright permissions}. APA Style Guidelines. \url{https://apastyle.apa.org/style-grammar-guidelines/citations/plagiarism}
\item Berne Convention. (1886, 9 de septiembre). \textit{Berne Convention for the Protection of Literary and Artistic Works} (París Act 1971, modificado en 1979). World Intellectual Property Organization. \url{https://www.wipo.int/wipolex/es/treaties/ShowDetails.jsp?treaty_id=15}
\item Brand, A., Allen, L., Altman, M., Hlava, M., \& Scott, J. (2015). Beyond authorship: Attribution, contribution, collaboration, and credit. \textit{Learned Publishing}, 28(2), 151--155. \url{https://doi.org/10.1087/20150211}
\item Bretag, T. (2013). Challenges in addressing student plagiarism: The perspectives of chief academic integrity officers. \textit{Ethics \& Behavior}, 23(4), 257--272. \url{https://doi.org/10.1080/10508422.2013.782006}
\item Bretag, T. (Ed.). (2016). \textit{Handbook of academic integrity}. Springer. \url{https://doi.org/10.1007/978-981-287-098-8}
\item Committee on Publication Ethics [COPE]. (2019). \textit{Core practices for journal editors and publishers}. COPE. \url{https://publicationethics.org/core-practices}
\item Committee on Publication Ethics [COPE]. (2023). \textit{Authorship and AI tools: COPE position statement}. COPE. \url{https://publicationethics.org/cope-position-statements/ai-author}
\item Corona Domínguez, B., \& Alonso Nuñez, M. C. (2026). Plagio académico y herramientas de inteligencia artificial generativa: Una revisión sistemática sobre implicaciones éticas y normativas en universidades mexicanas. \textit{Ciencia y Reflexión}, 5(1), 797--825. \url{https://doi.org/10.70747/cr.v5i1.926}
\item Dawson, P. (2021). \textit{Defending assessment security in a digital world: Preventing e-cheating and contract cheating in higher education}. Routledge. \url{https://doi.org/10.4324/9780429324178}
\item De Jesús G., M. I. (2016). Entre la ética en la investigación y la propiedad intelectual: Prácticas antiuniversitarias con relevancia para el derecho de autor. \textit{Actualidad Contable FACES}, 19(32), 40--67. \url{https://biblat.unam.mx/hevila/ActualidadcontableFACES/2016/vol19/no32/2.pdf}
\item Drahos, P. (2016). \textit{A philosophy of intellectual property}. Routledge. \url{https://doi.org/10.4324/9781315262796}
\item Eaton, S. E. (2021). \textit{Plagiarism in higher education: Tackling tough topics in academic integrity}. Libraries Unlimited. \url{https://doi.org/10.5040/9798400697333}
\item Eaton, S. E. (2023). Postplagiarism: Towards an era of artificial intelligence and academic integrity. \textit{Learning, Teaching and Research Forum}, 2(1), 1--5. \url{https://doi.org/10.11575/PRISM/40789}
\item European Commission. (2019). \textit{Ethics guidelines for trustworthy AI}. High-Level Expert Group on Artificial Intelligence. European Commission. \url{https://digital-strategy.ec.europa.eu/en/library/ethics-guidelines-trustworthy-ai}
\item European Commission. (2024). \textit{Living guidelines on the responsible use of generative AI in research}. Directorate-General for Research and Innovation. Publications Office of the European Union. \url{https://doi.org/10.2777/34947}
\item Foltýnek, T., Dlabolová, D., Qi, A., Janssen, M., \& Sivasubramaniam, S. (2020). Testing of support tools for plagiarism detection. \textit{International Journal of Educational Technology in Higher Education}, 17(1), 34. \url{https://doi.org/10.1186/s41239-020-00192-4}
\item Foltýnek, T., Bjelobaba, S., Glendinning, I., Khan, Z. R., Santos, R., Pavletic, P., \& Kravjar, J. (2023). ENAI recommendations on the ethical use of artificial intelligence in education. \textit{International Journal for Educational Integrity}, 19(1), 19. \url{https://doi.org/10.1007/s40979-023-00133-4}
\item Gallant, T. B. (2008). \textit{Academic integrity in the twenty-first century: A teaching and learning imperative}. ASHE Higher Education Report, 33(5). Jossey-Bass. \url{https://doi.org/10.1002/aehe.3305}
\item Goldstein, P., \& Hugenholtz, B. (2019). \textit{International copyright: Principles, law, and practice} (4th ed.). Oxford University Press. \url{https://doi.org/10.1093/oso/9780190060930.001.0001}
\item Hettinger, E. C. (1989). Justifying intellectual property. \textit{Philosophy \& Public Affairs}, 18(1), 31--52. \url{https://www.jstor.org/stable/2265419}
\item Howard, R. M. (1995). Plagiarisms, authorships, and the patchwriting toward of a pedagogy. \textit{College English}, 57(7), 788--806. \url{https://doi.org/10.2307/378403}
\item Howard, R. M. (1999). \textit{Standing in the shadow of giants: Plagiarists, authors, collaborators}. Ablex Publishing. ISBN: 978-1-56750-469-9.
\item Indecopi. (2018). \textit{Guía de derecho de autor para el ámbito universitario y de investigación}. Dirección de Derecho de Autor del Instituto Nacional de Defensa de la Competencia y de la Protección de la Propiedad Intelectual. \url{https://www.gob.pe/indecopi}
\item International Center for Academic Integrity. (2021). \textit{The fundamental values of academic integrity} (3.ª ed.). International Center for Academic Integrity. \url{https://www.academicintegrity.org/aws/ICAI/asset_manager/get_file/911282?ver=1}
\item International Committee of Medical Journal Editors. (2023). \textit{Recommendations for the conduct, reporting, editing, and publication of scholarly work in medical journals}. ICMJE. \url{http://www.icmje.org/icmje-recommendations.pdf}
\item Ji, Z., Lee, N., Frieske, R., Yu, T., Su, D., Yan, F., Ishii, E., Bang, Y. J., Dai, W., \& Fung, P. (2023). Survey of hallucination in natural language generation. \textit{ACM Computing Surveys}, 55(12), 248, 1--38. \url{https://doi.org/10.1145/3571730}
\item Landes, W. M., \& Posner, R. A. (2003). \textit{The economic structure of intellectual property law}. Harvard University Press. \url{https://doi.org/10.4159/9780674039919}
\item Lang, J. M. (2013). \textit{Cheating lessons: Learning from academic dishonesty}. Harvard University Press. \url{https://doi.org/10.4159/harvard.9780674726239}
\item Lessig, L. (2001). \textit{The future of ideas: The fate of the commons in a connected world}. Random House. ISBN: 978-0-375-50578-2.
\item Lessig, L. (2004). \textit{Free culture: How big media uses technology and the law to lock down culture and control creativity}. Penguin Press. ISBN: 978-1-59420-006-9.
\item Locke, J. (1690). \textit{Two treatises of government} (P. Laslett, Ed., 1988). Cambridge University Press. \url{https://doi.org/10.1017/CBO9780511810268}
\item Luenberger, D. G. (1997). \textit{Optimization by vector space methods}. John Wiley \& Sons. ISBN: 978-0-471-18117-0.
\item Macfarlane, B. (2017). The paradox of research integrity in a globalized knowledge economy. \textit{Educational Philosophy and Theory}, 49(3), 254--263. \url{https://doi.org/10.1080/00131857.2016.1174098}
\item Maynez, J., Narayan, S., Bohnet, B., \& McDonald, R. (2020). On faithfulness and factuality in abstractive summarization. \textit{Proceedings of the 58th Annual Meeting of the Association for Computational Linguistics (ACL)}, 1906--1919. \url{https://doi.org/10.18653/v1/2020.acl-main.173}
\item McCabe, D. L. (2005). Cheating among college and university students: A North American perspective. \textit{International Journal for Educational Integrity}, 1(1), 1--11. \url{https://doi.org/10.1007/s40979-005-0001-1}
\item McCabe, D. L., Butterfield, K. D., \& Treviño, L. K. (2012). \textit{Cheating in college: Why students do it and what educators can do about it}. Johns Hopkins University Press. \url{https://doi.org/10.1353/book.18431}
\item McCabe, D. L., Treviño, L. K., \& Butterfield, K. D. (2001). Cheating in academic institutions: A decade of research. \textit{Ethics \& Behavior}, 11(3), 219--232. \url{https://doi.org/10.1207/S15327019EB1103_2}
\item Merton, R. K. (1973). \textit{The sociology of science: Theoretical and empirical investigations}. University of Chicago Press. ISBN: 978-0-226-52092-6.
\item Miao, F., \& Holmes, W. (2023). \textit{Guidance for generative AI in education and research}. UNESCO Publishing. \url{https://unesdoc.unesco.org/ark:/48223/pf0000386693_eng}
\item Netanel, N. W. (2008). \textit{Copyright's paradox}. Oxford University Press. \url{https://doi.org/10.1093/acprof:oso/9780195137620.001.0001}
\item NISO. (2022). \textit{CRediT (Contributor Roles Taxonomy)} (ANSI/NISO Z39.104-2022). National Information Standards Organization. \url{https://www.niso.org/standards-committees/credit}
\item Ouyang, L., Wu, J., Jiang, X., Almeida, D., Wainwright, C. L., Mishkin, P., Zhang, C., Agarwal, S., Slama, K., Ray, A., Schulman, J., Hilton, J., Kelton, F., Miller, L., Simens, M., Askell, A., Welinder, P., Christiano, P., Leike, J., \& Lowe, R. (2022). Training language models to follow instructions with human feedback. \textit{Advances in Neural Information Processing Systems (NeurIPS)}, 35, 27730--27744. \url{https://doi.org/10.48550/arXiv.2203.02155}
\item Palacios Carrasco, R., Avalos Palomino, I., Noel Vargas, A. H., Cajas Aquino, L. A., \& Lavado Choque, A. X. (2026). \textit{Ética, plagio y propiedad intelectual en la universidad del siglo XXI: Navegando la integridad académica, los derechos de autor y el desafío de la Inteligencia Artificial Generativa en la investigación científica} [Monografía académica no publicada]. Escuela Profesional de Computación Científica, Facultad de Ciencias Matemáticas, Universidad Nacional Mayor de San Marcos.
\item Pecorari, D. (2003). Good and forbidden speech: Second language writers and the unauthorized use of source texts. \textit{English for Specific Purposes}, 22(4), 317--345. \url{https://doi.org/10.1016/S0889-4906(02)00025-9}
\item Pecorari, D. (2008). \textit{Academic writing and plagiarism: A linguistic analysis}. Continuum. ISBN: 978-0-8264-9246-3.
\item Perú. (1996, 24 de abril). \textit{Decreto Legislativo N.º 822: Ley sobre el Derecho de Autor}. Diario Oficial El Peruano. WIPO Lex. \url{https://www.wipo.int/wipolex/es/text/129301}
\item Perú. (2014, 9 de julio). \textit{Ley N.º 30220: Ley Universitaria}. Diario Oficial El Peruano. \url{https://www.gob.pe/institucion/sunedu/normas-legales/196621-30220}
\item Porto Castro, A. M. (2024). El plagio en la universidad: Una mirada pedagógica desde el rigor de los datos. En J. M. Touriñán López \& A. Rodríguez Martínez (Coords.), \textit{Pedagogía y educación, en perspectiva mesoaxiológica: Cuestiones conceptuales} (pp. 225--240). Editorial Redipe. \url{https://editorial.redipe.org/index.php/1/catalog/book/161}
\item Posner, R. A. (2007). \textit{The little book of plagiarism}. Pantheon Books. ISBN: 978-0-375-42475-5.
\item Puche Villalobos, D. J. (2025). La inteligencia artificial y el fraude académico en el contexto universitario. \textit{Revista Digital de Investigación y Postgrado}, 6(11), 73--93. \url{https://doi.org/10.59654/kg944e15}
\item Ricketson, S., \& Ginsburg, J. C. (2006). \textit{International copyright and neighbouring rights: The Berne Convention and beyond} (2nd ed.). Oxford University Press. \url{https://doi.org/10.1093/acprof:oso/9780198298717.001.0001}
\item Roig, M. (2001). Plagiarism and self-plagiarism: What every author should know. \textit{Croatian Medical Journal}, 42(4), 481--483.
\item Roig, M. (2015). \textit{Avoiding plagiarism, self-plagiarism, and other questionable writing practices: A guide to ethical writing}. Office of Research Integrity, U.S. Department of Health and Human Services. \url{https://ori.hhs.gov/avoiding-plagiarism-self-plagiarism-and-other-questionable-writing-practices-guide-ethical-writing}
\item Saldaña, J. (2021). La integridad científica en el sistema universitario peruano: Desafíos normativos e institucionales. \textit{Revista de Derecho y Sociedad}, 56, 115--138. \url{https://doi.org/10.18800/dys.202101.006}
\item Suber, P. (2012). \textit{Open access}. MIT Press. \url{https://doi.org/10.7551/mitpress/9286.001.0001}
\item Sutherland-Smith, W. (2008). \textit{Plagiarism, the Internet and student learning: Improving academic integrity}. Routledge. \url{https://doi.org/10.4324/9780203928370}
\item Swales, J. M., \& Feak, C. B. (2012). \textit{Academic writing for graduate students: Essential tasks and skills} (3rd ed.). University of Michigan Press. ISBN: 978-0-472-03475-8.
\item UNESCO. (2021). \textit{Recomendación sobre la ética de la inteligencia artificial}. Organización de las Naciones Unidas para la Educación, la Ciencia y la Cultura. \url{https://unesdoc.unesco.org/ark:/48223/pf0000381137_spa}
\item Vaidhyanathan, S. (2001). \textit{Copyrights and copywrongs: The rise of intellectual property and how it threatens creativity}. New York University Press. ISBN: 978-0-8147-8806-6.
\item Vaswani, A., Shazeer, N., Parmar, N., Uszkoreit, J., Jones, L., Gomez, A. N., Kaiser, \L., \& Polosukhin, I. (2017). Attention is all you need. \textit{Advances in Neural Information Processing Systems (NeurIPS)}, 30, 5998--6008. \url{https://doi.org/10.48550/arXiv.1706.03762}
\item Weber-Wulff, D., Anohina-Naumeca, A., Bjelobaba, S., Foltýnek, T., Guerrero-Dib, J., Popoola, O., Sigut, P., \& Waddington, L. (2023). Testing of detection tools for AI-generated text. \textit{International Journal for Educational Integrity}, 19(1), 26. \url{https://doi.org/10.1007/s40979-023-00146-z}
\item WIPO. (2016). \textit{Understanding copyright and related rights}. World Intellectual Property Organization. \url{https://www.wipo.int/publications/en/details.jsp?id=4081}
\item WIPO. (2023). \textit{Generative AI: Navigating intellectual property}. World Intellectual Property Organization. \url{https://www.wipo.int/about-ip/en/frontier_technologies/ai_and_ip.html}
\end{itemize}
